
\dictentry{Carbon Dioxide}{A colorless greenhouse gas composed of one carbon atom and two oxygen atoms, present in trace amounts in Earth's atmosphere but critical for regulating planetary temperature. In spacecraft life support systems, CO2 must be actively removed from cabin air using sorbent canisters or regenerable systems to protect crew health. On other planets like Venus, extremely high CO2 concentrations create a runaway greenhouse effect with surface temperatures exceeding 460 degrees Celsius.}{catSS}

\dictentry{Carbon Fiber Structure}{A spacecraft structural component fabricated from carbon fiber reinforced polymer composites that offer exceptional strength-to-weight ratios compared to traditional aluminum or titanium. Carbon fiber structures are used for bus panels, instrument booms, and antenna reflectors where minimizing mass while maintaining stiffness is critical. These composites require careful thermal design since their properties can degrade at very high temperatures encountered during launch or reentry.}{catSS}

\dictentry{Celestial Mechanics}{The branch of astronomy and physics that studies the motions of celestial bodies under the influence of gravitational forces. It provides the mathematical framework for predicting orbits, designing trajectories, and understanding phenomena such as tidal forces and orbital resonances. Classical celestial mechanics is based on Newtonian gravity, while modern applications incorporate relativistic corrections for high-precision orbit determination.}{catSS}

\dictentry{Center of Mass}{The theoretical point at which the entire mass of a spacecraft can be considered to be concentrated for the purposes of analyzing translational motion. It is calculated as the mass-weighted average of the positions of all components and must be maintained within tight tolerances for proper propulsion, attitude control, and thermal performance. Shifting center of mass due to propellant consumption requires careful management throughout a mission, and it is precisely measured during mass properties testing using specialized balance equipment on the ground before launch.}{catSS}

\dictentry{Ceres}{Ceres is the largest object in the asteroid belt between Mars and Jupiter, with a diameter of about 940 kilometers, and is classified as a dwarf planet by the International Astronomical Union. NASA's Dawn spacecraft orbited Ceres from 2015 to 2018, revealing bright salt deposits called faculae in Occator crater and evidence of subsurface briny water. Its relatively low density suggests a composition of rock and ice, making it a transitional body between the inner rocky planets and outer icy worlds.}{catSS}

\dictentry{Circular Velocity}{The specific velocity an object must maintain to stay in a perfectly circular orbit around a central body at a given altitude. For Earth orbit, circular velocity at the surface would be approximately 7.9 kilometers per second, decreasing with altitude. Any deviation from circular velocity will cause the orbit to become elliptical, with the object accelerating or decelerating along its path, and the resulting trajectory is described by Kepler's laws of planetary motion.}{catSS}

\dictentry{Clamp Band}{A tensioned metallic band that secures a payload to a launch vehicle adapter ring during ascent and releases on command to separate the payload in orbit. The band is pre-tensioned to a calculated force and released by a pyrotechnic or mechanical actuator that opens the band and allows springs to push the payload away. Clamp bands are widely used for their reliability, low shock, and simplicity in single-payload separations.}{catSS}

\dictentry{Clean Room}{A controlled-environment facility with filtered air, regulated temperature and humidity, and strict contamination control protocols used for assembling and testing spacecraft hardware. Particulate levels are continuously monitored and kept to extremely low counts to prevent contamination of sensitive optical instruments, solar cells, and electronic components. Personnel wear specialized garments including bunny suits, gloves, and face masks to minimize human-generated particles.}{catSS}

\dictentry{Cleanliness}{The controlled management of particulate and molecular contamination on spacecraft surfaces and components to prevent degradation of optical systems, thermal surfaces, and electronics. Cleanliness levels are specified in standards such as IEST-STD-CC1246, which defines cleanliness levels based on allowable particle sizes and quantities. Contamination can scatter light in telescopes, reduce solar cell efficiency, and cause electrical shorts in sensitive circuits.}{catSS}

\dictentry{Co-orbital}{Describing two or more objects that share the same orbit around a common body, potentially at different phases along that orbit. Co-orbital configurations include satellites in the same geostationary ring separated by longitude, or objects in Trojan arrangements sharing an orbit with a planet at Lagrange points. Co-orbital rendezvous and proximity operations require precise station-keeping to avoid collisions.}{catSS}

\dictentry{Cold Plate}{A heat-exchanger plate, typically made of aluminum or copper with internal fluid channels, used to remove heat from electronics and other heat-generating components in a spacecraft. Cold coolant or a phase-change fluid circulates through embedded tubes or channels, absorbing heat conducted from mounted components and transporting it to radiators. They are a key component of active thermal control systems where passive cooling alone is insufficient.}{catSS}

\dictentry{Comet}{A small icy body that releases gas and dust when heated by the Sun, developing a visible coma and sometimes tails that can extend millions of kilometers. Comets originate from the Kuiper Belt or Oort Cloud and are composed of frozen water, carbon dioxide, methane, and dust accumulated from the early solar system. Studying comets provides valuable information about the primordial material from which the planets formed over four billion years ago.}{catSS}

\dictentry{Compression}{The process of reducing the size of data before transmission from a spacecraft to ground stations to conserve bandwidth and power. Lossless compression algorithms such as Rice coding preserve all original data and are used for science telemetry, while lossy methods may be applied to imagery where some detail loss is acceptable. Effective compression ratios depend on the entropy of the source data and the algorithm employed.}{catSS}

\dictentry{Control Moment Gyro}{A momentum-exchange device consisting of a spinning rotor mounted in gimbal frames that generates torque by changing the angular momentum vector direction. Unlike reaction wheels, control moment gyros can produce much larger torques by tilting the spinning rotor, making them suitable for large spacecraft like the ISS. They allow continuous attitude slewing without consuming propellant, but must periodically dump accumulated momentum using external torques.}{catSS}

\dictentry{Convolutional Code}{An error-correcting code in which each output bit is a linear function of the current and several previous input bits, implemented using shift registers and modulo-2 adders. Convolutional codes provide good error correction performance with relatively simple encoding and decoding hardware, and are widely used in deep-space communication links. The Viterbi algorithm is commonly used for maximum-likelihood decoding of convolutional codes, and they are often concatenated with block codes for enhanced protection.}{catSS}

\dictentry{Coronal Mass Ejection}{A massive expulsion of magnetized plasma from the Sun's corona into interplanetary space, often associated with solar flares and filament eruptions. CMEs can carry billions of tonnes of material at speeds ranging from 250 to over 3,000 kilometers per second, and when directed toward Earth, can cause severe geomagnetic storms. These storms can disrupt satellite operations, power grids, and communication systems while producing spectacular auroral displays.}{catSS}

\dictentry{Cosmic Microwave Background}{The faint thermal radiation filling all of space, remnant heat from approximately 380,000 years after the Big Bang when the universe cooled enough for atoms to form and light to travel freely. Detected as a nearly uniform microwave signal at about 2.725 Kelvin, the CMB provides crucial evidence for the Big Bang theory and the large-scale structure of the universe. Minute fluctuations in the CMB reveal the seeds of galaxy formation and confirm models of cosmic inflation.}{catSS}

\dictentry{Cosmology}{The scientific study of the origin, evolution, large-scale structure, and ultimate fate of the universe as a whole. Cosmologists investigate questions about the Big Bang, dark matter, dark energy, cosmic expansion, and the fundamental laws governing spacetime. Modern cosmology combines theoretical physics, observational astronomy, and computational modeling to construct models such as the Lambda-CDM model that describe the universe's composition and history.}{catSS}


\dictentry{Cubesat}{A class of miniaturized satellite built to standardized dimensions of 10 by 10 by 10 centimeters per unit, with mass typically under 2 kilograms per unit. CubeSats use commercial off-the-shelf components to dramatically reduce development cost and time, making space accessible to universities and small organizations. Multiple units can be stacked and deployed from a single launch vehicle using standardized dispensers such as the Poly-Picosatellite Orbital Deployer.}{catSS}

\dictentry{Dark Energy}{A hypothetical form of energy permeating all of space that drives the observed accelerating expansion of the universe, comprising roughly 68 percent of the total energy content. Unlike dark matter, dark energy appears to have uniform density and exerts a repulsive gravitational effect. Its fundamental nature remains one of the deepest unsolved problems in physics, with competing theories including a cosmological constant and dynamic scalar fields.}{catSS}

\dictentry{Dark Matter}{A form of matter that does not emit, absorb, or reflect electromagnetic radiation but is detected through its gravitational effects on visible matter, radiation, and the large-scale structure of the universe. Dark matter comprises approximately 27 percent of the universe's total mass-energy content and is essential for explaining galaxy rotation curves, gravitational lensing, and cosmic structure formation. Its particle identity remains unknown, with leading candidates including weakly interacting massive particles.}{catSS}

\dictentry{Data Rate}{The speed at which data is transmitted between a spacecraft and ground stations, measured in bits per second. Data rate depends on transmitter power, antenna gain, distance to the receiver, and the modulation scheme used. Deep-space missions often operate at rates of only a few kilobits per second due to enormous distances, while low-Earth orbit satellites can achieve gigabit-per-second links using optical communication.}{catSS}

\dictentry{Dawn-Dusk Orbit}{A Sun-synchronous orbit in which the satellite passes over any given point on Earth at roughly the same local solar time near sunrise or sunset. The orbital plane is oriented approximately 90 degrees to the Sun direction, so the satellite rides along the terminator between day and night. This configuration provides near-continuous solar illumination of the spacecraft's solar panels while keeping instruments in favorable lighting conditions for Earth observation.}{catSS}

\dictentry{Deep Space Network}{NASA's global network of large antenna complexes located at Goldstone in California, Madrid in Spain, and Canberra in Australia, providing continuous communication coverage with spacecraft throughout the solar system. The three stations are spaced approximately 120 degrees apart in longitude so that as the Earth rotates, at least one station always has line-of-sight to any deep-space probe. The DSN uses antennas up to 70 meters in diameter and extremely sensitive receivers to detect faint signals from missions billions of kilometers away.}{catSS}

\dictentry{Deployable Structure}{Any spacecraft component designed to be stowed in a compact configuration during launch and mechanically unfurled or extended to its operational shape once in orbit. Deployable structures include solar arrays, antenna reflectors, booms, sunshields, and solar sails that must unfold reliably without ground intervention. Deployment mechanisms use stored energy in springs, motors, or shape-memory materials and are among the highest-risk events in a mission.}{catSS}

\dictentry{Descending Node}{The point in an orbit where the orbiting body crosses the reference plane, typically Earth's equatorial plane, moving from north to south. The descending node, together with the ascending node, defines the line of nodes and determines the orientation of the orbital plane in space. The positions of these nodes precess over time due to gravitational perturbations, particularly the J2 effect from Earth's equatorial bulge.}{catSS}

\dictentry{Despin}{The process of reducing or eliminating the spin rate of a spacecraft or stage after it has been spinning for stability during a previous mission phase. Despin can be achieved using thrusters, yo-yo weights, or by transferring angular momentum to a desaturation mechanism. Spin-down is necessary before deploying appendages, performing orbital maneuvers, or establishing three-axis stabilization for normal operations.}{catSS}

\dictentry{Docking Mechanism}{A mechanical system that enables two spacecraft to rendezvous, approach, and form a rigid, pressurized connection in orbit. Docking mechanisms include capture latches, structural hooks, seals, and alignment guides that progressively mate the two vehicles from initial contact through hard capture. International standards such as the International Docking System Standard ensure interoperability between spacecraft from different nations and manufacturers.}{catSS}

\dictentry{Drag Perturbation}{The gradual deceleration of a spacecraft in low Earth orbit caused by collisions with residual atmospheric molecules, which reduces orbital energy and causes the orbit to spiral inward. Atmospheric drag depends on the spacecraft's ballistic coefficient, cross-sectional area, and the local atmospheric density, which varies with solar activity and altitude. Without periodic reboosts, drag will eventually cause a satellite to reenter and burn up in the atmosphere.}{catSS}

\dictentry{DSN}{Deep Space Network, NASA's international array of large radio antennas used for communicating with and tracking interplanetary spacecraft. The network comprises three complexes at Goldstone, Madrid, and Canberra that together provide 360-degree coverage as the Earth rotates. The DSN also supports radio science experiments, very long baseline interferometry, and radar observations of asteroids and planetary bodies.}{catSS}

\dictentry{Dwarf Planet}{A celestial body massive enough to be rounded by its own gravity but that has not cleared its orbital neighborhood of other debris, as defined by the International Astronomical Union in 2006. Examples include Pluto, Eris, Haumea, Makemake, and Ceres, which orbit the Sun and may have moons but differ from full planets in orbital dominance. Dwarf planets are important targets for understanding the diversity of bodies in the outer solar system.}{catSS}

\dictentry{Earth Orbit}{Any orbital path around the Earth, classified by altitude into low Earth orbit below 2,000 kilometers, medium Earth orbit between 2,000 and 35,786 kilometers, and geostationary orbit at exactly 35,786 kilometers. Different orbits suit different missions: LEO for crewed flight and Earth observation, MEO for navigation satellites, and GEO for communications. The choice of orbit determines the satellite's velocity, period, coverage area, and communication requirements.}{catSS}

\dictentry{Earth Sensor}{An attitude determination sensor that detects the Earth's limb or infrared horizon to establish a local vertical reference for spacecraft orientation. Earth sensors typically use infrared detectors tuned to the 14 to 16 micrometer carbon dioxide absorption band to reliably identify the planet's edge regardless of lighting conditions. They are commonly used on low-Earth orbit and geostationary satellites for nadir pointing and attitude correction.}{catSS}

\dictentry{Eccentric Anomaly}{An angular parameter used in Keplerian orbital mechanics that relates the position of a body in an elliptical orbit to a circumscribed auxiliary circle. It serves as an intermediary variable in solving Kepler's equation to convert mean anomaly to true anomaly. Eccentric anomaly ranges from 0 to 360 degrees over one orbital period and provides a geometric construction for computing the satellite's position along its orbit, enabling efficient numerical orbit propagation.}{catSS}

\dictentry{Eccentricity}{A dimensionless parameter ranging from 0 to 1 that describes the shape of an orbit, where 0 represents a perfect circle and values approaching 1 represent increasingly elongated ellipses. Eccentricity is calculated from the ratio of the difference between apoapsis and periapsis distances to their sum. Orbits with eccentricity greater than 1 are hyperbolic, indicating the object is on an escape trajectory.}{catSS}

\dictentry{ECLSS}{Environmental Control and Life Support System, the integrated set of hardware that provides a habitable environment inside a crewed spacecraft or space station. ECLSS manages atmospheric composition, pressure, temperature, humidity, water recovery, and waste processing to sustain human life in the closed environment of space. The ISS ECLSS recycles approximately 90 percent of wastewater and continuously monitors and scrubs carbon dioxide from cabin air.}{catSS}

\dictentry{Electrodynamic Tether}{A long conductive cable deployed in orbit that generates electrical current through interaction with Earth's magnetic field as it cuts through magnetic field lines. The resulting Lorentz force can be used for propulsion or deorbiting without propellant, making electrodynamic tethers an attractive technology for station-keeping and space debris removal. The concept was demonstrated by the TSS-1R experiment on the Space Shuttle, though challenges remain with arcing and tether survivability.}{catSS}

\dictentry{EMI/EMC Test}{Electromagnetic interference and electromagnetic compatibility testing that verifies spacecraft electronics can operate correctly in their intended electromagnetic environment without causing or suffering harmful interference. Tests include conducted and radiated emissions measurements, susceptibility testing using RF fields and electrical fast transients, and electrostatic discharge testing. Compliance with EMC standards is mandatory before flight to ensure subsystems coexist without performance degradation.}{catSS}

\dictentry{Engineering Model}{A non-flight representation of spacecraft hardware built to the same design as the flight article but using commercial-grade components, used during the detailed design phase for functional testing and interface verification. Engineering models allow engineers to validate software, test electrical interfaces, and refine procedures before committing to expensive flight hardware. They are typically maintained throughout the program for troubleshooting and regression testing.}{catSS}

\dictentry{Error Correction}{The technique of adding redundant bits to transmitted data so that errors introduced during communication can be detected and corrected at the receiver without retransmission. Space communications use error correction extensively because signals from deep space arrive with very low signal-to-noise ratios and cannot be easily retransmitted. Common schemes include convolutional codes, Reed-Solomon codes, turbo codes, and LDPC codes, often concatenated for maximum reliability.}{catSS}

\dictentry{Escape Velocity}{The minimum speed an object must achieve to break free from the gravitational attraction of a body without further propulsion. For Earth, escape velocity from the surface is approximately 11.2 kilometers per second, while for the Moon it is only 2.4 kilometers per second. Escape velocity depends on the mass and radius of the central body and determines whether an object follows a bound orbit or an escape trajectory.}{catSS}

\dictentry{EVA Suit}{An Extravehicular Activity suit, a self-contained spacecraft that provides life support, thermal regulation, micrometeorite protection, and mobility for astronauts working outside a pressurized vehicle. The suit maintains a pressurized atmosphere, supplies oxygen, removes carbon dioxide, and shields against temperatures ranging from minus 150 to plus 120 degrees Celsius. Modern EVA suits such as the EMU used on the ISS allow several hours of independent work in the vacuum of space.}{catSS}

\dictentry{Exoplanet}{A planet that orbits a star outside our solar system, detected through methods such as radial velocity measurements, transit photometry, or direct imaging. Thousands of exoplanets have been discovered since the first confirmed detection in 1992, revealing a remarkable diversity of worlds including super-Earths, hot Jupiters, and planets in their star's habitable zone. The study of exoplanets is central to understanding planet formation and the potential for life elsewhere in the universe.}{catSS}

\dictentry{Exosphere}{The outermost and thinnest layer of a planet's atmosphere, extending from the top of the thermosphere to the point where particles escape into space. In Earth's exosphere, starting at about 500 kilometers altitude, gas molecules are so sparse that collisions between them are rare and some can follow ballistic trajectories or escape entirely. This region contains the upper portions of the Van Allen radiation belts and is where many low-Earth orbit satellites operate.}{catSS}

\dictentry{Flight Model}{The actual flight-ready spacecraft or hardware unit that will be launched into space, built and tested to the highest quality standards with flight-qualified components. The flight model undergoes a rigorous test campaign including vibration, acoustic, thermal vacuum, and electromagnetic compatibility testing to ensure it can survive the launch environment and operate in space. Only one flight model is typically produced per mission, making it extremely valuable.}{catSS}

\dictentry{Flyby Spacecraft}{A spacecraft that follows a trajectory passing close to a celestial body without entering orbit, using the brief encounter to gather scientific data and take images. Flyby missions such as Voyager, New Horizons, and Galileo's Jupiter tour cover vast distances and can visit multiple targets using gravity assists. The limited encounter time, often only hours or days of close approach, requires careful planning to maximize data return from each target.}{catSS}

\dictentry{Frangible Nut}{A pyrotechnic separation device consisting of a nut with an internally housed explosive charge that shatters the nut when initiated, instantly releasing a bolted joint. Frangible nuts are used throughout spacecraft and launch vehicles to release hold-down mechanisms, stage separation joints, and deployment latches with high reliability and minimal shock. Once fired, the nut fragments are captured by a containment sleeve to prevent debris generation.}{catSS}

\dictentry{Free Fall}{The motion of an object subject solely to gravitational force with no other forces acting on it, resulting in apparent weightlessness for objects in orbit. In low Earth orbit, spacecraft and everything inside them are in continuous free fall around the Earth, producing the microgravity environment experienced by ISS crew. True free fall is rarely achieved perfectly, as residual atmospheric drag and tidal forces create small accelerations that define the microgravity quality of the orbiting laboratory.}{catSS}

\dictentry{Frozen Orbit}{A specific orbit where perturbation effects from the oblateness of the central body are balanced so that the orbit's eccentricity and argument of perigee remain essentially constant over time. For Earth, frozen orbits are typically near-circular with inclinations near 63.4 or 116.6 degrees where the J2 and J3 perturbation effects cancel. Frozen orbits are highly desirable for Earth observation missions that require consistent altitude over specific latitudes.}{catSS}

\dictentry{Fuel Cell}{An electrochemical device that continuously generates electricity by combining hydrogen fuel and oxygen from the air, producing water as a byproduct. Fuel cells provide higher energy density than batteries for long-duration missions and were used extensively on the Apollo spacecraft and Space Shuttle to generate electrical power while producing drinking water for the crew. NASA continues to develop advanced fuel cells for future crewed exploration missions.}{catSS}

\dictentry{Galaxy}{A massive gravitationally bound system containing hundreds of billions of stars along with stellar remnants, interstellar gas, dust, and dark matter. Galaxies range in size from dwarf galaxies with a few billion stars to giant ellipticals containing over a hundred trillion stars, and are classified by shape into spiral, elliptical, and irregular types. Our Milky Way is a barred spiral galaxy approximately 100,000 light-years in diameter containing 100 to 400 billion stars.}{catSS}

\dictentry{Gas Giant}{A large planet composed predominantly of hydrogen and helium with a relatively small rocky core, such as Jupiter and Saturn in our solar system. Gas giants have deep atmospheres with complex weather systems, strong magnetic fields, and numerous moons and ring systems. They play a crucial role in shaping the architecture of planetary systems by influencing the orbits of smaller bodies through their powerful gravitational fields, and studying them reveals insights into planetary formation processes.}{catSS}

\dictentry{Geostationary Orbit}{A circular equatorial orbit at an altitude of exactly 35,786 kilometers where a satellite's orbital period matches Earth's rotation period, making it appear stationary over a fixed point on the ground. This orbit is ideal for communications, weather, and broadcasting satellites that need continuous coverage of a specific region. Satellites in GEO must maintain precise station-keeping to counteract gravitational perturbations from the Sun and Moon.}{catSS}

\dictentry{Geostationary Transfer Orbit}{An elliptical orbit with perigee in low Earth orbit and apogee at geostationary altitude, used as an intermediate trajectory to deliver satellites to geostationary orbit. A spacecraft is first launched into GTO, then performs one or more circularization burns at apogee to raise perigee and reduce inclination until it reaches GEO. GTO is the most common delivery orbit for commercial communications satellites launched from equatorial or near-equatorial sites.}{catSS}

\dictentry{Geosynchronous Orbit}{Any orbit with an orbital period equal to Earth's sidereal rotation period of approximately 23 hours and 56 minutes, meaning the satellite returns to the same position relative to the ground each day. Unlike geostationary orbit, a geosynchronous orbit may be inclined, causing the satellite to trace a figure-eight ground track. This orbit class includes both true GEO and inclined orbits used for communications and navigation in high-latitude regions.}{catSS}

\dictentry{Graveyard Orbit}{A disposal orbit located several hundred kilometers above or below geostationary orbit where decommissioned GEO satellites are moved at end of life to prevent interference with operational spacecraft. International guidelines recommend a super-synchronous disposal orbit at least 235 kilometers above GEO to ensure the defunct satellite remains clear of the active orbital ring for at least 100 years. The disposal maneuver requires sufficient remaining fuel to perform the raising burn.}{catSS}

\dictentry{Gravitational Constant}{The fundamental physical constant G, approximately 6.674 times 10 to the minus 11 newton square meters per kilogram squared, that appears in Newton's law of universal gravitation. It determines the strength of the gravitational force between any two masses separated by a given distance. The value of G is known to less than one percent accuracy and is one of the least precisely measured fundamental constants in physics.}{catSS}

\dictentry{Gravity}{The universal force of attraction between all objects with mass, described by Newton's law of universal gravitation and more precisely by Einstein's general theory of relativity as the curvature of spacetime. Gravity governs the orbits of planets and satellites, shapes large-scale cosmic structure, and determines the internal forces in stars and planets. Despite being the weakest of the four fundamental forces, gravity is always attractive and acts over infinite range, making it dominant on cosmic scales.}{catSS}

\dictentry{Gravity Assist}{A spaceflight maneuver in which a spacecraft uses the relative motion and gravity of a planet or other celestial body to change its speed and direction without expending propellant. The spacecraft gains or loses energy relative to the Sun as it trades momentum with the planet during a carefully planned flyby. Gravity assists have enabled missions like Voyager, Cassini, and New Horizons to reach distant targets that would otherwise require impractically large rockets.}{catSS}

\dictentry{Gravity Gradient Stabilization}{A passive attitude control method that exploits the difference in gravitational force across the length of a long boom or extended structure to keep one axis of a satellite aligned with the local vertical. The near end of the satellite experiences slightly stronger gravity than the far end, creating a restoring torque that aligns the spacecraft's minimum moment-of-inertia axis toward the planet. This technique is simple and requires no propellant but is limited to low-Earth orbit applications.}{catSS}

\dictentry{Greenhouse Effect}{The process by which certain gases in a planet's atmosphere absorb and re-emit infrared radiation from the surface, trapping heat and raising the planet's average temperature above what it would be without an atmosphere. On Earth, water vapor, carbon dioxide, and methane are the primary greenhouse gases and are essential for maintaining habitable temperatures. Runaway greenhouse effects on Venus demonstrate how excessive greenhouse gas concentrations can make a planet uninhabitable.}{catSS}

\dictentry{Ground Station}{A terrestrial facility equipped with antennas, receivers, transmitters, and computers for communicating with and tracking spacecraft in orbit. Ground stations send commands to spacecraft and receive telemetry and science data downlinks, forming the critical link between mission control and the space segment. Major ground station networks include NASA's Near Space Network, the ESA ESTRACK system, and commercial providers such as KSAT and SSC.}{catSS}

\dictentry{Gyroscope}{A sensor that measures angular rate or rotation by exploiting the conservation of angular momentum of a spinning mass, or in modern MEMS and fiber-optic versions, through interference effects. Gyroscopes are essential components of spacecraft inertial measurement units, providing the rate data needed for attitude determination and control. Ring laser and fiber-optic gyroscopes offer high accuracy with no moving parts, while MEMS gyroscopes provide low-cost solutions for small satellites.}{catSS}

\dictentry{Habitat Module}{A pressurized module on a space station or spacecraft designed to provide living quarters, workspace, and life support for crew during long-duration missions. Habitat modules include sleeping areas, exercise equipment, galley facilities, and waste management systems, all within a controlled atmospheric environment. Examples include the ISS Unity and Zvezda modules and the planned Gateway Habitation and Logistics Outpost for lunar missions, which will test deep-space habitation technologies.}{catSS}

\dictentry{Halo Orbit}{A three-dimensional periodic orbit around a Lagrange point in which the spacecraft traces a roughly circular path out of the plane of the two primary bodies. Halo orbits require regular station-keeping maneuvers to maintain but offer continuous line-of-sight to both the Sun and a planet, making them useful for space observatories. The James Webb Space Telescope operates in a halo orbit around the Sun-Earth L2 point approximately 1.5 million kilometers from Earth.}{catSS}

\dictentry{Hard Capture}{The final, rigid, and fully sealed connection in a docking sequence where the two spacecraft are mechanically locked together and the docking interface is pressurized to allow crew and cargo transfer. Hard capture follows soft capture and involves extending structural hooks or bolts to create a structurally sound joint that can withstand attitude control forces. Once hard capture is complete, hatches can be opened between the two vehicles.}{catSS}

\dictentry{Heat Pipe}{A passive two-phase heat transfer device that uses the evaporation and condensation of a working fluid to transfer heat efficiently over relatively long distances with minimal temperature drop. The working fluid evaporates at the hot end, travels as vapor to the cold end where it condenses releasing latent heat, and returns to the hot end via capillary action through a wick structure. Heat pipes are widely used in spacecraft thermal control for spreading heat from electronics to radiator panels.}{catSS}

\dictentry{Heliosphere}{The vast region of space surrounding the Sun that is filled with the solar wind and the magnetic field carried outward by the solar wind, extending well beyond the planets. The heliosphere shields the solar system from a significant portion of incoming galactic cosmic rays and interstellar particles. The Voyager spacecraft have crossed the heliopause, the boundary where the solar wind meets the interstellar medium, entering interstellar space for the first time.}{catSS}

\dictentry{High Gain Antenna}{A directional antenna with a narrow beamwidth and high gain that concentrates transmitted power into a focused beam, enabling high data-rate communication over vast distances. High-gain antennas on spacecraft typically use parabolic reflectors or phased arrays and must be precisely pointed toward the receiving station. They are essential for deep-space missions where signal strength falls off with the square of the distance, but cannot be used when the spacecraft's attitude is not precisely known.}{catSS}

\dictentry{Highly Elliptical Orbit}{An orbit with high eccentricity, featuring a low perigee and a very high apoapsis that causes the satellite to spend most of its orbital period near apoapsis where it moves slowly. This property is exploited by Molniya and Tundra orbits to provide prolonged communication coverage over high-latitude regions where geostationary satellites have poor visibility. Satellites in HEO can remain visible to a ground station for many hours per orbit.}{catSS}

\dictentry{Hill Sphere}{The region around a celestial body within which it can gravitationally retain satellites despite the tidal pull of a more massive body. For Earth, the Hill sphere extends about 1.5 million kilometers, roughly the distance to the L1 Lagrange point. Moons and satellites must orbit within their primary's Hill sphere to maintain a stable orbit over long periods, which is why the Moon orbits Earth rather than the Sun directly.}{catSS}

\dictentry{Hohmann Transfer Orbit}{The most fuel-efficient two-impulse maneuver for transferring a spacecraft between two coplanar circular orbits, consisting of an elliptical transfer orbit tangent to both the initial and target orbits. The first burn raises the orbit from the initial circle, and the second burn circularizes at the target altitude. Although fuel-efficient, Hohmann transfers require a specific transfer angle and travel time, which may not always match mission timing requirements.}{catSS}

\dictentry{Honeycomb Panel}{A lightweight structural panel consisting of a honeycomb core sandwiched between two thin face sheets, used extensively in spacecraft structures for its excellent stiffness-to-weight ratio. The honeycomb core, typically made of aluminum or Nomex, provides high shear strength and bending stiffness while the face sheets carry tension and compression loads. These panels are used for satellite bus structures, instrument mounting platforms, and launch vehicle payload adapters.}{catSS}

\dictentry{Horizon Sensor}{An attitude determination instrument that detects the infrared radiation from a planet's atmospheric limb to establish a local vertical reference for nadir pointing. The sensor measures the contrast between the warm planet and cold space to determine the spacecraft's roll and pitch angles relative to the horizon. Horizon sensors are commonly used on geostationary communication satellites and Earth observation platforms for maintaining precise ground-pointing accuracy.}{catSS}

\dictentry{Impact Crater}{A circular depression on the surface of a celestial body created by the high-velocity collision of a meteoroid, asteroid, or comet with the surface. Impact craters range from microscopic pits on lunar samples to the 1,800-kilometer-wide Hellas Basin on Mars, and their size and distribution reveal the body's geological history and surface age. Craters are the primary mechanism for resurfacing airless bodies and serve as natural windows into subsurface geology.}{catSS}

\dictentry{Inclination}{The angle between a satellite's orbital plane and the reference plane, typically Earth's equatorial plane for Earth-orbiting satellites. Inclination ranges from 0 degrees for an equatorial orbit to 90 degrees for a polar orbit, and determines the latitudinal range over which the satellite passes directly overhead. Changing inclination is one of the most energetically expensive orbital maneuvers, requiring large delta-v burns.}{catSS}

\dictentry{Inertial Measurement Unit}{An integrated sensor package combining three orthogonal accelerometers and three orthogonal gyroscopes to measure a spacecraft's linear acceleration and rotational rate in six degrees of freedom. The IMU provides the raw data needed for inertial navigation, attitude determination, and guidance computations during all mission phases. Modern spacecraft IMUs use ring laser or fiber-optic gyroscopes for high accuracy and reliability over long-duration missions.}{catSS}

\dictentry{Integration}{The process of assembling, connecting, and verifying the interfaces between all spacecraft subsystems to form a complete, functional vehicle. Integration includes mechanical assembly, electrical harness routing, fluid line connections, and software loading, followed by comprehensive functional testing. This phase typically follows individual subsystem testing and is one of the most labor-intensive and critical phases of spacecraft development.}{catSS}

\dictentry{Integration Facility}{A specialized building equipped with clean rooms, overhead cranes, test equipment, and environmental controls where spacecraft and launch vehicle components are assembled and tested. The facility provides controlled conditions for mating the spacecraft with the launch vehicle adapter, conducting final inspections, and performing system-level verification tests. Major integration facilities include NASA's Vehicle Assembly Building and SpaceX's Horizontal Integration Facility.}{catSS}

\dictentry{Interplanetary Transfer}{A trajectory designed to move a spacecraft from one planet's gravitational sphere of influence to another's, following a heliocentric orbit that connects the departure and arrival orbits. Transfer trajectories are computed using patched conics approximation and are optimized for minimum fuel, minimum time, or specific encounter conditions. Gravity assists from intermediate planets are often incorporated to reduce the required launch energy.}{catSS}

\dictentry{Ionosphere}{The electrically ionized layer of Earth's upper atmosphere, extending from about 60 to 1,000 kilometers altitude, where ultraviolet and X-ray radiation from the Sun strips electrons from atoms and molecules. The ionosphere affects radio wave propagation by reflecting, refracting, and absorbing signals depending on their frequency, making it important for communication and navigation systems. It contains several distinct layers that change in density and altitude with time of day, season, and solar activity.}{catSS}

\dictentry{ISS}{The International Space Station, a permanently crewed modular laboratory in low Earth orbit at approximately 408 kilometers altitude, built and operated as a collaboration among NASA, Roscosmos, ESA, JAXA, and CSA. The ISS serves as a microgravity and space environment research laboratory where scientific experiments are conducted in astrobiology, astronomy, meteorology, and physics. It has been continuously inhabited since November 2000 and is the largest human-made object in low Earth orbit.}{catSS}

\dictentry{J2 Perturbation}{The dominant orbital perturbation caused by the oblateness of a planet, specifically the J2 zonal harmonic coefficient that represents the equatorial bulge resulting from the planet's rotation. J2 causes secular changes including precession of the ascending node, rotation of the line of apsides, and periodic variations in orbital elements. For Earth, J2 equals approximately 0.00108263 and must be accounted for in all precise orbit propagation and mission planning.}{catSS}

\dictentry{Ka Band}{A portion of the electromagnetic spectrum in the microwave frequency range from approximately 26.5 to 40 gigahertz, used for high-bandwidth space-to-ground communication links. Ka-band offers significantly higher data rates than lower-frequency bands such as S and X band, but is more susceptible to rain fade attenuation in the atmosphere. Modern high-throughput communication satellites and deep-space missions increasingly use Ka-band to meet growing data return requirements.}{catSS}

\dictentry{Kennedy Space Center}{NASA's primary launch complex on Merritt Island, Florida, adjacent to Cape Canaveral Space Force Station, serving as the agency's main center for human spaceflight operations. KSC includes the historic Launch Complex 39 where Apollo and Space Shuttle missions departed, as well as the Vehicle Assembly Building and Operations and Checkout Building. Today it supports commercial crew and cargo missions to the ISS and is being developed for the Artemis program returning humans to the Moon.}{catSS}

\dictentry{Kepler's First Law}{The first of Kepler's three laws of planetary motion, stating that the orbit of every planet is an ellipse with the Sun at one of the two foci. This law replaced the long-held belief that planets move in perfect circles, fundamentally changing our understanding of celestial mechanics. For spacecraft, this law applies equally well to any two-body gravitational system and is the basis for computing elliptical transfer orbits.}{catSS}

\dictentry{Kepler's Laws}{The three empirical laws formulated by Johannes Kepler in the early 17th century that describe the motion of planets around the Sun. They state that orbits are ellipses, that planets sweep equal areas in equal times, and that the square of the orbital period is proportional to the cube of the semi-major axis. Kepler's laws provided the foundation for Newton's theory of gravitation and remain fundamental to modern orbital mechanics and spacecraft trajectory design.}{catSS}

\dictentry{Kepler's Second Law}{The second of Kepler's three laws, stating that a line segment joining a planet and the Sun sweeps out equal areas during equal intervals of time. This means a planet moves faster when it is closer to the Sun at perihelion and slower when farther away at aphelion. The law is a direct consequence of the conservation of angular momentum and applies to all orbiting bodies including spacecraft and satellites.}{catSS}

\dictentry{Kepler's Third Law}{The third of Kepler's three laws, stating that the square of the orbital period of a planet is directly proportional to the cube of the semi-major axis of its orbit. This mathematical relationship allows calculation of orbital period from semi-major axis, or vice versa, and shows that more distant orbits have much longer periods. The law provides the basis for determining the mass of celestial bodies from their orbital characteristics.}{catSS}

\dictentry{Kourou}{The Guiana Space Centre, a European spaceport located near Kourou in French Guiana on the northeastern coast of South America. Its proximity to the equator at approximately 5 degrees north latitude provides a significant launch advantage, allowing rockets to gain extra velocity from Earth's rotation for geostationary orbit insertion. Kourou is the primary launch site for the European Ariane and Vega rocket families, supporting commercial, scientific, and institutional missions.}{catSS}

\dictentry{Kuiper Belt}{A vast region of the outer solar system extending from roughly 30 to 50 astronomical units from the Sun, beyond the orbit of Neptune, containing hundreds of thousands of icy bodies larger than 100 kilometers. The Kuiper Belt is home to several dwarf planets including Pluto, Haumea, and Makemake, and is thought to be the source of many short-period comets. It provides critical clues about the early formation and dynamical evolution of the outer solar system.}{catSS}

\dictentry{L1}{The first Lagrange point, located approximately 1.5 million kilometers from Earth on the line connecting the Sun and Earth, where the combined gravitational pull of both bodies balances the centrifugal force. A spacecraft at L1 can maintain a continuous view of the Sun-facing side of Earth and is ideal for solar observation missions such as SOHO and the DSCOVR space weather satellite. Station-keeping at L1 requires only small periodic corrections to maintain position.}{catSS}

\dictentry{L2}{The second Lagrange point, located approximately 1.5 million kilometers from Earth on the far side of the Sun-Earth line, where gravitational and centrifugal forces balance for a co-orbiting object. L2 is the preferred location for space telescopes that need to observe the cold, dark universe away from the Sun, including the James Webb Space Telescope, Planck, and Gaia. The stable thermal environment and continuous deep-space communication access make L2 ideal for infrared astronomy.}{catSS}

\dictentry{Lagrange Point}{One of five positions in a two-body orbital system where a small object can maintain a stable or quasi-stable position relative to the two larger bodies. At these points, the combined gravitational forces of the two bodies provide exactly the centripetal force needed to orbit with them. Lagrange points are widely used for space observatories, solar wind monitoring, and communications relay satellites because they offer unique vantage points with minimal fuel for station-keeping.}{catSS}

\dictentry{Lander}{A spacecraft designed to descend from orbit or flyby trajectory and make a controlled landing on the surface of a celestial body such as a planet, moon, or asteroid. Landers must withstand the stresses of atmospheric entry if applicable, perform precision descent using thrusters or parachutes, and survive touchdown with shock-absorbing legs or airbags. Successful landers such as Mars Pathfinder, Phoenix, and Chang'e-5 carry instruments to study surface composition, weather, and geology.}{catSS}

\dictentry{Launch Adapter}{A structural interface that mechanically and electrically connects a spacecraft to the upper stage of its launch vehicle, providing load paths for launch forces and separation mechanisms for deployment. The adapter typically includes ordnance for pyrotechnic separation, electrical connectors for pre-launch monitoring, and alignment features for proper spacecraft orientation. Its design must accommodate the specific interfaces of both the spacecraft and the launch vehicle while minimizing mass.}{catSS}

\dictentry{Launch Campaign}{The sequence of activities at a launch site beginning with the arrival of the launch vehicle and spacecraft and culminating in a successful liftoff. Activities include vehicle assembly, spacecraft integration, propellant loading, system checks, rehearsal countdowns, and the final launch sequence itself. The campaign typically lasts several weeks and involves hundreds of engineers, technicians, and mission controllers working in tight coordination.}{catSS}

\dictentry{Launch Complex}{A ground-based facility designed for preparing, fueling, and launching rockets, consisting of a launch pad, service structures, propellant storage, tracking systems, and mission control centers. Launch complexes are built to withstand the intense heat, vibration, and acoustic forces generated during liftoff while providing safe access for personnel during pre-launch operations. Each complex is typically tailored to specific launch vehicle families.}{catSS}



\dictentry{Launch Site}{A geographic facility chosen for rocket launches based on factors including latitude, range safety, downrange tracking, weather patterns, and logistics access. Optimal launch sites near the equator provide a velocity boost from Earth's rotation for eastward launches, while polar launch sites require careful avoidance of populated downrange areas. Major launch sites include Cape Canaveral, Baikonur, Kourou, Tanegashima, and Vandenberg.}{catSS}

\dictentry{Launch Window}{The specific period of time during which a spacecraft must be launched to reach its intended target orbit or celestial body, determined by orbital mechanics, target position, and lighting conditions. Launch windows may last from minutes to hours and are calculated months or years in advance based on the target body's position relative to Earth. Missing a launch window may require waiting weeks, months, or even years for the next favorable opportunity.}{catSS}

\dictentry{LDPC}{Low-Density Parity-Check codes, a class of linear error-correcting codes defined by sparse parity-check matrices that can achieve performance very close to the Shannon limit. LDPC codes are decoded using iterative belief propagation algorithms that efficiently correct errors even at very low signal-to-noise ratios. They are used in modern deep-space communication standards and are increasingly adopted for near-Earth high-data-rate links due to their excellent performance and parallelizable decoding.}{catSS}

\dictentry{Life Support System}{An integrated system of hardware that provides all environmental conditions necessary for human survival in space, including breathable atmosphere, potable water, temperature and humidity control, and waste management. Advanced life support systems close material loops by recycling water from humidity and urine, regenerating oxygen from carbon dioxide, and filtering contaminants from cabin air. Future systems for long-duration missions aim for near-complete closure to reduce resupply requirements.}{catSS}

\dictentry{Light-Year}{The distance that light travels in a vacuum in one Julian year, approximately 9.461 trillion kilometers or 63,241 astronomical units. Light-years are used to express the enormous distances between stars and galaxies, with the nearest star Proxima Centauri located about 4.24 light-years from Earth. Despite the name, a light-year is a unit of distance, not time, though it also represents how far back in time we observe an object.}{catSS}

\dictentry{Lightband}{A low-shock separation system developed by Planetary Systems Corporation that uses a motorized rotating band to release a payload from a launch vehicle adapter without pyrotechnics. The system gradually unwinds a titanium band that holds the payload, providing extremely gentle and controlled separation with minimal shock and debris. Lightband is used for satellite deployments and interstage separations where minimizing disturbance is critical for sensitive payloads.}{catSS}

\dictentry{Line of Nodes}{The straight line connecting the ascending node and descending node of an orbit, representing the intersection of the orbital plane with the reference plane such as Earth's equatorial plane. The orientation of the line of nodes is measured by the right ascension of the ascending node and changes over time due to nodal precession caused by gravitational perturbations. Understanding the line of nodes is essential for planning orbit-raising maneuvers and rendezvous operations.}{catSS}

\dictentry{Lissajous Orbit}{A quasi-periodic trajectory around a Lagrange point in which the spacecraft oscillates in both the in-plane and out-of-plane directions with different frequencies, tracing a complex non-repeating pattern. Unlike a halo orbit, a Lissajous orbit does not close on itself and can drift over time, requiring periodic station-keeping maneuvers. This orbit type is useful for missions that do not need continuous line-of-sight to both primary bodies, offering flexibility in fuel management.}{catSS}

\dictentry{Louver}{A variable-position thermal control device consisting of movable slats or vanes mounted over a radiator surface that can be opened or closed to regulate the rate of heat rejection to space. When louvers are open, the radiator can dissipate maximum heat; when closed, they insulate the radiator and retain heat within the spacecraft. Louvers provide an active yet mechanically simple method of thermal control that can respond to changing heat loads during different mission phases, and are commonly used on Earth-orbiting satellites with cyclic thermal environments.}{catSS}

\dictentry{Low Earth Orbit}{An orbit around Earth with an altitude between approximately 160 and 2,000 kilometers, where most crewed spaceflight and many Earth observation satellites operate. LEO offers the advantage of lower launch energy and higher data rates due to proximity, but satellites experience atmospheric drag that gradually decays their orbit. Orbital periods in LEO range from about 88 to 127 minutes, and the environment includes radiation belt passages and space debris.}{catSS}

\dictentry{Low Gain Antenna}{An omnidirectional or wide-beam antenna on a spacecraft that provides communication coverage over a broad angular range without requiring precise pointing. Low-gain antennas offer lower data rates than high-gain antennas but are essential for maintaining contact during attitude uncertainties, safe-mode operations, or early post-launch phases before the high-gain antenna is deployed and pointed. They serve as a reliable backup communication path throughout a mission.}{catSS}

\dictentry{Magnetic Torquer}{A device consisting of electromagnets or permanent magnets that interacts with a planet's magnetic field to generate torque for spacecraft attitude control. By selectively energizing coils aligned with different axes, the control system can apply torques to rotate the spacecraft without consuming propellant. Magnetic torquers are commonly used on low-Earth orbit satellites for momentum desaturation and coarse attitude adjustment, though their effectiveness decreases at higher altitudes where the magnetic field is weaker.}{catSS}

\dictentry{Magnetometer}{A scientific instrument that measures the strength and direction of magnetic fields in space, used both for science investigations and as an attitude determination sensor. Fluxgate and cesium vapor magnetometers mounted on long booms minimize interference from the spacecraft's own magnetic field. Magnetometers on planetary missions have mapped the magnetic fields of Jupiter, Saturn, and Mercury, revealing details about their interior structures and magnetospheric dynamics.}{catSS}

\dictentry{Magnetorquer}{An electromagnetic coil that generates a magnetic dipole moment when electric current flows through it, creating torque against a planetary magnetic field. Magnetorquers are lightweight, reliable actuators used extensively on small satellites and cubesats for attitude control and momentum dumping. They are often used in combination with reaction wheels, where the magnetorquers periodically unload accumulated angular momentum without propellant expenditure, extending mission life significantly.}{catSS}

\dictentry{Magnetosphere}{The region of space around a magnetized planet where the planet's magnetic field dominates over the solar wind, creating a protective bubble that deflects most charged particles. Earth's magnetosphere extends tens of thousands of kilometers on the sunward side and stretches into a long tail on the nightside. The magnetosphere contains the Van Allen radiation belts and is responsible for channeling charged particles toward the poles, producing auroras.}{catSS}

\dictentry{Margin of Safety}{The ratio by which a structure's design strength exceeds the maximum expected load, used as a design factor to account for uncertainties in material properties, manufacturing tolerances, load predictions, and analysis methods. Spacecraft structural margins of safety are typically required to be positive by a specified amount throughout the design, test, and flight loads environment. Negative margins indicate a design deficiency that must be resolved before flight, driving redesign or reinforcement of the affected structure.}{catSS}

\dictentry{Marman Clamp}{A cylindrical clamp band separation system that uses a V-shaped groove ring to hold two structural rings together under tension, commonly used for spacecraft-to-rocket-adapter joints and stage separations. When the tensioning bolts are released by pyrotechnic cutters, the band expands and allows springs to push the two bodies apart. Marman clamps provide reliable, well-characterized separation forces and are one of the most widely used mechanical separation systems in the launch industry.}{catSS}

\dictentry{Mass Properties Measurement}{The experimental determination of a spacecraft's total mass, center of gravity location, and moments and products of inertia using precision measurement equipment. These measurements are critical for verifying that the spacecraft meets its design mass budget, that the center of gravity falls within the launch vehicle adapter limits, and that attitude control system actuators can handle the inertia. Measurements are typically performed on a three-axis balance or pivot system.}{catSS}

\dictentry{Mast}{A slender, deployable vertical structure on a spacecraft used to position antennas, sensors, solar arrays, or other equipment at a distance from the main body. Masts are typically stowed as a compact bundle during launch and extended using spring-loaded hinges, telescoping mechanisms, or motorized drives. They provide structural support and electrical routing for deployed elements while minimizing mass and stowed volume, and are commonly fabricated from lightweight composites or aluminum alloys.}{catSS}

\dictentry{Mean Anomaly}{An angular parameter that specifies the position of a body along its orbit as a fraction of the total orbital period, assuming uniform circular motion. Unlike true anomaly, which varies non-uniformly with time in an elliptical orbit, mean anomaly increases at a constant rate and is directly proportional to elapsed time since periapsis passage. Converting mean anomaly to true anomaly requires solving Kepler's equation iteratively, a fundamental computation in all orbit determination and prediction software.}{catSS}

\dictentry{Mechanism}{Any moving or deployable mechanical component on a spacecraft, including hinges, latches, motors, gears, sliding joints, and release devices. Space mechanisms must operate reliably in the harsh environment of vacuum, extreme temperatures, and radiation, often after long storage periods on the ground. They are among the highest-risk elements of a spacecraft because failure of a single mechanism can compromise an entire mission, demanding rigorous testing and qualification.}{catSS}

\dictentry{Medium Earth Orbit}{An orbital altitude range between approximately 2,000 and 35,786 kilometers, encompassing the orbits used by navigation satellite constellations such as GPS at about 20,200 kilometers, GLONASS, Galileo, and BeiDou. MEO offers a balance between the low latency of low Earth orbit and the wide coverage of geostationary orbit, making it ideal for global positioning and navigation services. Satellites in MEO experience higher radiation exposure from the Van Allen belts than those in LEO.}{catSS}

\dictentry{Mesosphere}{The layer of Earth's atmosphere extending from about 50 to 85 kilometers altitude, above the stratosphere and below the thermosphere. In the mesosphere, temperature decreases with altitude due to the absence of ozone absorption, reaching the coldest temperatures in Earth's atmosphere at the mesopause. This layer is where most meteors burn up due to friction with atmospheric gases, producing the visible streaks we call shooting stars.}{catSS}

\dictentry{Meteor}{The visible streak of light produced when a meteoroid enters Earth's atmosphere at high speed and heats the surrounding air through compression, causing ionization and luminosity. Most meteors are caused by particles no larger than a grain of sand and burn up completely at altitudes of 70 to 100 kilometers. During meteor showers, dozens to hundreds of meteors per hour can be seen as Earth passes through the debris trail left by a comet.}{catSS}

\dictentry{Meteorite}{A fragment of rock or metal from space that survives passage through Earth's atmosphere and reaches the surface, providing direct samples of asteroids and other celestial bodies. Meteorites are classified into three main types: stony (chondrites and achondrites), iron, and stony-iron, each revealing different aspects of solar system history. They are invaluable for studying the age, composition, and processes of planetary formation because they preserve records dating back 4.5 billion years.}{catSS}

\dictentry{Meteoroid}{A small rocky or metallic body moving through interplanetary space, ranging in size from dust grains to objects about one meter across, smaller than asteroids. Meteoroids originate from collisions between asteroids, cometary debris trails, and ejecta from planetary surfaces. When a meteoroid enters Earth's atmosphere it becomes a meteor, and if any portion survives to the ground it is classified as a meteorite.}{catSS}

\dictentry{Microgravity}{The condition of near-weightlessness experienced by objects in free fall around Earth, where the only forces acting are gravitational, producing apparent gravitational acceleration of less than one millionth of Earth's surface gravity. Microgravity on the ISS is caused by the continuous free fall of the station around Earth, not by the absence of gravity. This environment enables unique scientific research in fluid physics, materials science, combustion, and human physiology that cannot be replicated on the ground.}{catSS}

\dictentry{Milky Way}{The barred spiral galaxy containing our Solar System, with a diameter of approximately 100,000 to 180,000 light-years containing an estimated 100 to 400 billion stars. The Milky Way has a central bar structure surrounded by spiral arms rich in young stars, gas, and dust, with a supermassive black hole called Sagittarius A* at its center. From Earth, the galaxy appears as a luminous band across the night sky, with darker lanes caused by interstellar dust obscuring more distant stars.}{catSS}

\dictentry{Molniya Orbit}{A highly elliptical orbit with an inclination of 63.4 degrees, an orbital period of approximately 12 hours, and a high apogee over the northern hemisphere, named after the Russian communications satellites that first used it. The special inclination freezes the argument of perigee, keeping apogee fixed at high latitudes where geostationary satellites provide poor coverage. Molniya orbits provide extended dwell times over northern Russia and other high-latitude regions.}{catSS}

\dictentry{Moon}{A natural satellite orbiting a planet, formed through various processes including giant impact, gravitational capture, or co-accretion from a protoplanetary disk. Earth's Moon, the fifth largest satellite in the solar system, stabilizes Earth's axial tilt, generates tides, and provides a relatively accessible destination for human exploration. Many moons in the outer solar system, such as Europa and Enceladus, are prime candidates in the search for extraterrestrial life due to their subsurface oceans.}{catSS}

\dictentry{Multi-Layer Insulation}{A thermal control blanket consisting of multiple thin layers of aluminized Mylar or Kapton separated by low-conductivity spacers, used to minimize heat transfer by radiation on spacecraft. MLI is effective in vacuum because it eliminates conductive heat transfer between layers and the reflective surfaces minimize radiative exchange. The distinctive gold or silver appearance of many satellites comes from their MLI blankets, which are carefully tailored to each spacecraft's thermal requirements.}{catSS}

\dictentry{Nebula}{A vast cloud of interstellar gas and dust that may span hundreds of light-years and serves as the birthplace of new stars or the remnant of dead ones. Nebulae are classified into emission nebulae lit by hot young stars, reflection nebulae that scatter starlight, dark nebulae that obscure background stars, and planetary nebulae expelled by dying stars. The Orion Nebula and Crab Nebula are among the most studied examples, revealing details about stellar birth and death.}{catSS}

\dictentry{Neutron Star}{An extraordinarily dense stellar remnant formed when a massive star collapses during a supernova explosion, packing roughly 1.4 times the Sun's mass into a sphere only about 20 kilometers in diameter. The matter is so compressed that protons and electrons merge into neutrons, and the surface gravity is hundreds of billions of times Earth's. Rapidly spinning neutron stars that emit beams of radiation are called pulsars, and those in binary systems can accrete matter and emit X-rays.}{catSS}

\dictentry{Newton's Law of Gravitation}{The fundamental law stating that every particle of matter in the universe attracts every other particle with a force proportional to the product of their masses and inversely proportional to the square of the distance between their centers. Formulated by Isaac Newton in 1687, it provides accurate predictions for orbital mechanics and trajectory design for spacecraft throughout the solar system. While superseded by general relativity for extreme conditions, Newton's law remains sufficient for virtually all engineering applications.}{catSS}

\dictentry{Nodal Precession}{The rotation of an orbit's line of nodes around the central body over time, caused primarily by the oblateness of the planet through the J2 perturbation effect. For low-Earth orbits, nodal precession can cause the orbital plane to rotate several degrees per day, which is exploited in Sun-synchronous orbits to maintain a constant angle with the Sun. The rate and direction of nodal precession depend on the orbit's altitude, inclination, and eccentricity.}{catSS}

\dictentry{Northern Lights}{The Aurora Borealis, a natural luminous display in Earth's high-latitude sky caused by charged particles from the solar wind being channeled along magnetic field lines into the upper atmosphere. The particles collide with nitrogen and oxygen molecules, exciting them to higher energy states that produce characteristic green, red, purple, and blue emissions as they return to ground state. The intensity and frequency of auroras are closely tied to the solar cycle and geomagnetic storm activity.}{catSS}

\dictentry{Nuclear Power in Space}{The use of nuclear fission reactors or radioisotope decay to generate electrical power and heat for spacecraft operating far from the Sun or in shadowed environments where solar panels are ineffective. Radioisotope thermoelectric generators convert heat from plutonium-238 decay into electricity and have powered missions like Voyager, Cassini, and Curiosity. Kilopower fission reactors under development could provide tens of kilowatts for future lunar bases and Mars surface operations.}{catSS}

\dictentry{Oort Cloud}{A theoretical spherical shell of icy objects extending from roughly 2,000 to 100,000 astronomical units from the Sun, believed to be the source of long-period comets that occasionally enter the inner solar system. The Oort Cloud has never been directly observed but is inferred from the orbits of comets and the dynamical history of the solar system. Its existence suggests that material from the early solar system was scattered outward by gravitational interactions with the giant planets.}{catSS}

\dictentry{Orbital Mechanics}{The application of Newtonian mechanics and gravitational theory to predict and control the motion of objects in space, encompassing orbit determination, propagation, and maneuver planning. It provides the mathematical tools for designing transfer orbits, calculating delta-v budgets, predicting ground tracks, and planning rendezvous and proximity operations. Orbital mechanics is essential for every aspect of spaceflight from launch vehicle trajectory design to satellite constellation management.}{catSS}

\dictentry{Orbital Period}{The time required for a satellite or celestial body to complete one full orbit around the central body, determined by the semi-major axis and the mass of the central body through Kepler's third law. For low Earth orbit the period is about 90 minutes, for geosynchronous orbit it is approximately 24 hours, and for the Moon it is about 27.3 days. The orbital period determines how frequently a satellite passes over a given location on the surface.}{catSS}

\dictentry{Orbital Velocity}{The speed a spacecraft must maintain along its orbital path to remain in orbit around a central body, balancing gravitational attraction with the centrifugal effect of its motion. Orbital velocity decreases with altitude, from about 7.8 kilometers per second in low Earth orbit to 3.1 kilometers per second at geostationary altitude. The velocity also varies within an elliptical orbit, being highest at periapsis and lowest at apoapsis.}{catSS}

\dictentry{Orbiter}{A spacecraft designed to enter and remain in orbit around a celestial body for extended periods, conducting scientific observations, communications relay, or reconnaissance of the surface below. Orbiters carry suites of remote sensing instruments including cameras, spectrometers, and radar to study the body's atmosphere, surface, gravity field, and magnetic environment from orbit. Successful orbiters include Mars Reconnaissance Orbiter, Cassini, and Juno.}{catSS}

\dictentry{Outgassing}{The release of volatile gases and trapped molecules from materials when exposed to the vacuum of space, caused by the evaporation of adsorbed moisture, solvents, and low-molecular-weight compounds. Outgassing can deposit films on sensitive optical surfaces, contaminate thermal control coatings, and degrade material properties over time. Space-qualified materials are tested for outgassing rates using standards such as ASTM E595 to ensure compatibility with the space environment.}{catSS}

\dictentry{Ozone Layer}{A region of Earth's stratosphere between approximately 15 and 35 kilometers altitude containing a high concentration of ozone molecules that absorb harmful ultraviolet-B and ultraviolet-C radiation from the Sun. The ozone layer is essential for protecting life on Earth from DNA-damaging UV radiation, and its depletion by chlorofluorocarbons was a major environmental concern addressed by the Montreal Protocol. Recovery of the ozone layer demonstrates the effectiveness of international environmental policy.}{catSS}

\dictentry{Parabolic Antenna}{A directional antenna using a parabolic reflector surface to focus incoming radio waves onto a feed horn or to collimate transmitted signals into a narrow beam, providing high gain and directivity. Parabolic antennas on spacecraft and ground stations achieve very narrow beamwidths suitable for high-data-rate deep-space communication and radar applications. The antenna's gain increases with the ratio of the reflector diameter to the wavelength of the signal.}{catSS}

\dictentry{Parking Orbit}{A temporary orbit where a spacecraft resides after launch and before performing a subsequent maneuver to reach its final operational orbit or interplanetary trajectory. Parking orbits allow ground controllers to perform system checks, wait for proper phasing, or conduct orbit-raising burns at optimal points. Most geostationary satellite missions use a parking orbit in low Earth orbit before performing the transfer to GEO.}{catSS}

\dictentry{Parsec}{A unit of distance used in astronomy equal to approximately 3.26 light-years or 30.86 trillion kilometers, defined as the distance at which one astronomical unit subtends an angle of one arcsecond. Parsecs are preferred by astronomers over light-years because they relate directly to the trigonometric parallax method of measuring stellar distances. One kiloparsec equals 1,000 parsecs and is used to measure distances within our galaxy.}{catSS}

\dictentry{Passive Thermal Control}{A spacecraft thermal management approach using fixed elements such as multi-layer insulation, surface coatings, thermal paints, heat sinks, and heat pipes to maintain component temperatures within acceptable ranges without active intervention. Passive systems are favored for their simplicity, reliability, and zero power consumption. However, they have limited ability to respond to changing thermal loads and are typically combined with active elements for missions with diverse thermal requirements.}{catSS}

\dictentry{Patched Conics}{An approximation method for computing interplanetary trajectories by dividing the overall path into segments where only one gravitational body dominates, with each segment described by a conic section. Within a planet's sphere of influence the trajectory is a hyperbolic escape or capture curve, and between spheres it is a heliocentric ellipse. This method provides useful first-order trajectory designs that can be refined with higher-fidelity numerical integration.}{catSS}

\dictentry{Payload Adapter}{A structural interface that connects a satellite or spacecraft payload to the launch vehicle's upper stage, transmitting launch loads while providing a separation interface for deployment in orbit. The adapter must match the mechanical, electrical, and thermal interfaces of both the payload and the rocket, and typically includes ordnance for pyrotechnic separation. Its design is critical for ensuring clean, shock-minimized payload deployment.}{catSS}

\dictentry{Payload Processing}{The series of activities performed on a spacecraft or satellite at the launch site to prepare it for flight, including final inspections, fueling, battery charging, configuration loading, and integration with the launch vehicle. Payload processing occurs in specialized facilities with clean room environments and specialized equipment. The processed payload is then transported to the integration facility for mating with the launch vehicle.}{catSS}

\dictentry{Periapsis}{The point in an orbit closest to the central body, where the orbiting object reaches its maximum velocity and experiences the strongest gravitational pull. The periapsis altitude and the apoapsis together define the eccentricity and size of the orbit. At periapsis, spacecraft may experience the most intense atmospheric drag, heating, or tidal forces depending on the mission profile, and maneuvers executed at this point are often the most efficient for changing orbital energy.}{catSS}

\dictentry{Perigee}{The point in an Earth orbit closest to the Earth's center, where the satellite reaches its maximum orbital velocity. Perigee altitude is one of the two parameters, along with apogee, that define an elliptical Earth orbit. Low perigee values result in increased atmospheric drag that can accelerate orbital decay, while high-altitude perigees are characteristic of transfer orbits such as GTO that are used to reach geostationary altitude.}{catSS}

\dictentry{Perihelion}{The point in a solar orbit where a body is closest to the Sun, where it moves at its maximum orbital velocity as described by Kepler's second law. Earth reaches perihelion around early January each year, at a distance of approximately 147.1 million kilometers. For comets, the intense solar heating at perihelion causes their volatile ices to sublimate, creating the visible coma and tails.}{catSS}

\dictentry{Perturbation}{A deviation of an orbiting body's trajectory from a perfect two-body Keplerian orbit, caused by additional forces such as atmospheric drag, solar radiation pressure, gravitational influence of other bodies, or the oblateness of the central body. Perturbations cause gradual changes in orbital elements including precession of nodes, rotation of the line of apsides, and changes in eccentricity and semi-major axis. Accurate orbit propagation requires modeling these perturbative effects.}{catSS}

\dictentry{Phase Change Material}{A substance that absorbs or releases large amounts of latent heat when it transitions between solid and liquid states, used in spacecraft thermal control to buffer temperature fluctuations. During heating the material melts, absorbing energy without a temperature rise, and during cooling it solidifies, releasing stored heat. Phase change materials are particularly useful for instruments and electronics that experience cyclic heat loads during eclipse and sunlit periods.}{catSS}

\dictentry{Phased Array}{An antenna system in which multiple radiating elements are combined with electronically controlled phase shifters to steer the antenna beam without physically moving the antenna structure. Phased arrays can redirect their beam in microseconds, enabling rapid tracking of ground stations or targets and supporting multiple simultaneous communication links. They are increasingly used on modern communication and radar satellites for their flexibility and reliability compared to mechanically steered dishes.}{catSS}

\dictentry{Phasing Orbit}{An orbit used to adjust the relative position of a spacecraft with respect to a target, typically by changing the orbital period so that the spacecraft arrives at a specific point at a desired time. Phasing maneuvers involve short burns that slightly change the semi-major axis, allowing the spacecraft to gain or lose angular position relative to the target over subsequent orbits. Phasing orbits are essential for rendezvous operations and constellation maintenance.}{catSS}

\dictentry{Photovoltaic Cell}{A semiconductor device that directly converts photons from sunlight into electrical energy through the photovoltaic effect, forming the basis of solar panels and arrays on spacecraft. Space-grade solar cells are typically made from gallium arsenide or multi-junction compounds to maximize efficiency under the intense solar radiation and radiation environment of space. Conversion efficiencies for modern triple-junction cells exceed 30 percent under standard test conditions.}{catSS}

\dictentry{Pin Puller}{A pyrotechnic actuator that retracts a retaining pin to release a spring-loaded mechanism or latch, commonly used to initiate deployment sequences or stage separations on spacecraft. When fired by an electrical command, a small explosive charge drives a piston that extracts the pin, allowing the restrained mechanism to activate. Pin pullers are valued for their simplicity, reliability, and clean operation without generating significant debris or shock.}{catSS}

\dictentry{Planet}{A celestial body massive enough to achieve hydrostatic equilibrium, resulting in a roughly spherical shape, that orbits a star and has cleared its orbital neighborhood of other debris. The eight planets of our solar system range from small rocky worlds like Mercury to gas giants like Jupiter, each with distinct atmospheres, surface conditions, and satellite systems. Exoplanet discoveries have revealed thousands of additional planets orbiting other stars in a remarkable variety of configurations.}{catSS}

\dictentry{Planetary Protection}{The set of guidelines and practices designed to prevent biological contamination of celestial bodies by terrestrial organisms and to protect Earth from potential extraterrestrial life forms. Implemented through NASA's Office of Planetary Protection, these measures include strict clean room assembly requirements, sterilization of spacecraft destined for biologically sensitive targets like Mars, and containment protocols for sample return missions. Planetary protection is essential for preserving the scientific integrity of life-detection experiments.}{catSS}

\dictentry{Planetary Science}{The scientific study of planets, moons, dwarf planets, asteroids, comets, and other bodies in the solar system and beyond, combining geology, atmospheric science, oceanography, and astrophysics. Planetary scientists analyze data from spacecraft missions, telescopic observations, meteorite studies, and laboratory experiments to understand the formation, evolution, and current conditions of worlds throughout the universe. This field has been revolutionized by missions such as Voyager, Mars rovers, and the James Webb Space Telescope.}{catSS}

\dictentry{Pluto}{A dwarf planet in the Kuiper Belt at approximately 39 astronomical units from the Sun, with a diameter of about 2,377 kilometers and five known moons including the large moon Charon. Pluto's surface features nitrogen and methane ices, water ice mountains, and a heart-shaped plain called Tombaugh Regio, revealed by NASA's New Horizons flyby in 2015. Its highly elliptical and inclined orbit periodically brings it closer to the Sun than Neptune.}{catSS}

\dictentry{Polar Orbit}{An orbit with an inclination near 90 degrees that passes over or near both poles of the Earth on each revolution, providing complete global surface coverage over time. Polar orbits at altitudes of 700 to 800 kilometers are commonly used for Earth observation, reconnaissance, and weather monitoring satellites because they allow the satellite to scan every longitude as the Earth rotates beneath. The high inclination also provides good coverage of polar regions that geostationary satellites cannot observe.}{catSS}

\dictentry{Poynting-Robertson Effect}{A drag force on small dust particles in orbit around the Sun caused by the asymmetry in absorbed and re-emitted solar radiation due to the particle's orbital motion. The particle absorbs sunlight in its forward direction of motion but re-emits isotropically, creating a net deceleration that slowly spirals the particle inward toward the Sun. This effect is significant for micron-sized dust grains and shapes the distribution of interplanetary dust in the inner solar system, and also affects the orbital evolution of small near-Earth asteroids over long timescales.}{catSS}

\dictentry{Probe and Drogue}{A two-component docking system in which an active spacecraft extends a probe into a conical receptacle called a drogue on the passive spacecraft, providing initial alignment and capture before hard mate. The probe tip engages the drogue cone and latches, then retracts to pull the two vehicles together for structural connection. This system was used on Soyuz, Apollo, and early ISS docking ports, though newer systems favor androgynous designs.}{catSS}

\dictentry{Prograde Orbit}{An orbit in which the satellite moves in the same direction as the central body's rotation, resulting in a positive inclination measured from the equatorial plane. Prograde orbits benefit from the rotational velocity boost of the Earth, requiring less launch energy than retrograde orbits for the same altitude. Most operational satellites including LEO Earth observation and GEO communications satellites use prograde orbits.}{catSS}

\dictentry{Propellant Loading}{The process of transferring cryogenic and storable propellants from ground storage facilities into the launch vehicle's tanks in preparation for launch. Cryogenic propellants such as liquid oxygen and liquid hydrogen require insulated transfer lines and continuous boil-off management, while storable propellants like hypergolics are handled at ambient temperature. Propellant loading is one of the most hazardous operations during the countdown and is closely monitored for leaks and tank pressurization.}{catSS}

\dictentry{Protoflight Model}{A single spacecraft or hardware unit that serves dual roles as both the qualification test article and the flight unit, undergoing a reduced-level test campaign to verify design margins while preserving the hardware for launch. This approach saves cost and schedule compared to building separate qualification and flight models, but accepts some risk because the hardware is stressed during testing. Protoflight testing typically uses lower test levels than full qualification to maintain sufficient margin.}{catSS}

\dictentry{Pyrotechnic Device}{An explosive actuator that converts stored chemical energy into rapid mechanical motion for functions such as stage separation, antenna deployment, fairing jettison, and valve actuation on spacecraft and launch vehicles. Pyrotechnic devices include bolts, nuts, cutters, linear shaped charges, and thrusters that provide extremely reliable, fast-acting, and compact mechanisms. They are single-use devices that must function perfectly on the first command after potentially years of storage.}{catSS}

\dictentry{Qualification Model}{A complete or representative spacecraft hardware unit built specifically for design validation testing, subjecting it to environmental stresses beyond expected flight levels to demonstrate adequate design margin. The qualification model is tested to higher levels than the flight article, typically 125 percent of predicted maximum environments, to prove the design is robust. Qualification test results are used to certify the design for flight without risking the actual flight hardware.}{catSS}

\dictentry{Radiation Pressure}{A small but persistent force exerted on a spacecraft by photons from sunlight or other electromagnetic radiation, caused by the transfer of momentum from absorbed or reflected photons. Though tiny at about 9 micronewtons per square meter at Earth's distance from the Sun, radiation pressure causes significant orbital perturbations over long periods and is the basis for solar sail propulsion. It must be accounted for in precise orbit determination and station-keeping calculations.}{catSS}

\dictentry{Radiator}{A specialized surface on a spacecraft designed to reject excess heat to space by emitting infrared thermal radiation, typically coated with high-emissivity paint or surfaces. Radiators are connected to the spacecraft's thermal control system via heat pipes or fluid loops and must be oriented to minimize solar heating while maximizing view to deep space. The sizing of radiators is driven by the worst-case heat dissipation and the allowable temperature range of the components they cool.}{catSS}

\dictentry{Radioisotope Thermoelectric Generator}{A power source that converts heat from the natural decay of radioactive isotopes, typically plutonium-238, into electricity using thermoelectric couples. RTGs produce no moving parts and require no maintenance, making them ideal for deep-space missions where solar power is impractical, such as Voyager, Cassini, and the New Horizons mission to Pluto. A single RTG can produce hundreds of watts of electrical power for decades, though output gradually declines as the isotope decays.}{catSS}



\dictentry{Reaction Wheel}{A momentum-storage device consisting of a spinning flywheel mounted on a motorized axis within the spacecraft, used to exchange angular momentum for attitude control without consuming propellant. By changing the spin rate of one or more reaction wheels, the spacecraft rotates in the opposite direction while conserving total angular momentum. Sets of four or more wheels arranged in skewed configurations provide three-axis control with built-in redundancy.}{catSS}

\dictentry{Redshift}{The increase in wavelength of electromagnetic radiation from an astronomical object, caused either by the object moving away from the observer through the Doppler effect or by the expansion of space itself stretching the wavelength during propagation. Cosmological redshift is evidence for the expansion of the universe, with more distant galaxies showing greater redshifts. The redshift of spectral lines allows astronomers to measure distances, velocities, and the chemical composition of distant objects.}{catSS}

\dictentry{Reed-Solomon Code}{A non-binary cyclic error-correcting code that operates on blocks of data symbols rather than individual bits, providing excellent burst-error correction capability for space communications. Reed-Solomon codes can correct multiple symbol errors within a block and are particularly effective at handling the bursty errors common in wireless channels. They were used extensively in deep-space missions and are often concatenated with convolutional codes for a powerful two-layer error protection scheme.}{catSS}

\dictentry{Regolith}{The layer of loose, fragmented, and unconsolidated material covering the surface of solid celestial bodies, formed by billions of years of meteorite impacts, volcanic activity, and space weathering. Lunar regolith is a fine-grained mixture of glass spherules, mineral fragments, and metallic iron particles ranging from dust to boulders, and presents challenges for spacecraft landing and surface operations. Understanding regolith properties is crucial for planning future lunar and Martian surface missions and resource utilization.}{catSS}

\dictentry{Remote Manipulator System}{A robotic arm system mounted on a spacecraft or space station for capturing, deploying, and manipulating payloads, scientific instruments, and other objects in space. The Canadarm2 on the ISS is a 17.6-meter-long RMS with seven degrees of freedom that can grapple payloads up to 116 tonnes and is essential for station assembly, maintenance, and cargo vehicle capture. The arm provides both teleoperated and autonomous capabilities for complex orbital operations.}{catSS}

\dictentry{Retrograde Orbit}{An orbit in which the satellite moves in the direction opposite to the central body's rotation, resulting in an inclination greater than 90 degrees. Retrograde orbits require significantly more launch energy than prograde orbits because the rocket must overcome rather than benefit from Earth's rotational velocity. While rare for Earth-orbiting satellites, retrograde orbits are common for solar observation missions that require specific sun angles.}{catSS}

\dictentry{RF Communication}{Radio frequency communication using electromagnetic waves in the radio spectrum for transmitting commands, receiving telemetry, and relaying science data between spacecraft and ground stations. RF links must account for free-space path loss, Doppler shift, signal polarization, and atmospheric effects, and employ modulation and coding schemes optimized for the available power and bandwidth. Modern spacecraft may use S, X, Ka, or laser bands depending on data rate requirements and mission constraints.}{catSS}

\dictentry{Right Ascension of Ascending Node}{The angle measured in the equatorial plane from the vernal equinox direction to the ascending node, specifying the orientation of the orbital plane in inertial space. RAAN is one of the six classical Keplerian orbital elements and is typically expressed in hours, minutes, and seconds of right ascension. Changes in RAAN caused by J2 perturbation are exploited to maintain Sun-synchronous orbits where the RAAN precesses at the same rate as the Sun's apparent motion.}{catSS}

\dictentry{Robotic Arm}{A multi-jointed manipulator mechanism on a spacecraft or space station that can be controlled by astronauts or ground operators to grasp, move, and release objects in the microgravity environment. Robotic arms range from simple two-joint mechanisms to complex seven-degree-of-freedom systems with force-torque sensors and cameras for precise proximity operations. They are essential for station assembly, satellite servicing, and capture of free-flying cargo vehicles.}{catSS}

\dictentry{Roche Limit}{The minimum distance from a celestial body at which a second body held together only by its own gravity can orbit without being torn apart by tidal forces from the primary. Inside the Roche limit, tidal forces exceed the self-gravity holding the smaller body together, causing it to disintegrate into a ring or debris field. Saturn's rings are thought to have formed from a moon or captured body that was disrupted after passing within the planet's Roche limit.}{catSS}

\dictentry{Rover}{A mobile robotic vehicle designed to traverse and explore the surface of a celestial body, carrying instruments for geological analysis, atmospheric sampling, and imaging. Mars rovers such as Spirit, Opportunity, Curiosity, and Perseverance have traveled tens of kilometers across Martian terrain, drilling rock samples and searching for signs of past water and biological activity. Rovers must navigate autonomously or semi-autonomously due to communication delays with Earth.}{catSS}

\dictentry{RTG}{Radioisotope Thermoelectric Generator, a long-lasting power source for deep-space missions that converts heat from radioactive decay into electricity. RTGs use thermoelectric couples to generate power from the temperature difference between the hot isotope heat source and the cold radiators facing deep space. A typical RTG produces between 100 and 300 watts of electrical power and can operate reliably for decades with no moving parts or maintenance.}{catSS}

\dictentry{S Band}{A portion of the electromagnetic spectrum in the frequency range from approximately 2 to 4 gigahertz, widely used for spacecraft telemetry, tracking, and command links. S-band offers a good balance between data rate capability and resistance to atmospheric attenuation, making it suitable for both near-Earth and deep-space communications. It is the standard frequency band for launch vehicle tracking and for initial communication with spacecraft after separation from the launch vehicle.}{catSS}

\dictentry{Sample Return Mission}{A space mission designed to collect physical samples from a celestial body such as an asteroid, comet, or planet and return them to Earth for detailed laboratory analysis. Sample return missions involve complex operations including surface collection, ascent from the body, rendezvous in space, and safe Earth reentry and recovery. Missions such as OSIRIS-REx, Hayabusa2, and Chang'e-5 have demonstrated increasingly sophisticated sample return techniques.}{catSS}

\dictentry{Satellite}{An artificial object placed into orbit around a celestial body to serve functions including communications, Earth observation, navigation, weather monitoring, scientific research, and military surveillance. Satellites range in size from tiny cubesats weighing a few kilograms to massive geostationary platforms exceeding six tonnes. They operate in a variety of orbits selected to optimize coverage, data relay, or observational requirements for their specific mission.}{catSS}

\dictentry{Semi-Major Axis}{Half of the longest diameter of an elliptical orbit, representing the average distance of the orbiting body from the central body and serving as a primary measure of the orbit's size. The semi-major axis directly determines the orbital period through Kepler's third law, with larger axes corresponding to longer periods. It is one of the six classical Keplerian orbital elements that uniquely describe a satellite's orbit.}{catSS}

\dictentry{Semi-Minor Axis}{Half of the shortest diameter of an elliptical orbit, perpendicular to the semi-major axis at the center of the ellipse and related to the semi-major axis through the eccentricity. Together with the semi-major axis and eccentricity, it fully defines the shape and size of the orbital ellipse. The semi-minor axis is always less than or equal to the semi-major axis, reaching equality only for a perfectly circular orbit where the eccentricity is zero.}{catSS}

\dictentry{Separation Nut}{A pyrotechnic fastener consisting of a nut with an internal explosive charge that, when initiated, fractures the nut body along pre-scored lines to release a bolted joint. Separation nuts are used for stage separation, payload deployment, and mechanism release where rapid, reliable, and clean disconnection is required. Once fired, the separated halves of the nut are restrained by the bolt to prevent free-floating debris, and the design ensures no loose fragments are released into space.}{catSS}

\dictentry{Separation System}{The complete set of mechanical, pyrotechnic, and electrical components that physically disconnect the spacecraft from the launch vehicle or from a previous stage at the designated point in the mission timeline. Separation systems must provide clean, symmetric push-off forces, minimal shock, and reliable function after long storage. Common types include clamp bands, marman clamps, spring pushers, and linear shaped charges.}{catSS}

\dictentry{Shock Test}{A test that exposes spacecraft hardware to the high-amplitude, short-duration mechanical impulses generated by pyrotechnic events such as bolt releases, stage separations, and deployment mechanisms. Pyrotechnic shock contains energy at very high frequencies up to hundreds of kilohertz and can damage sensitive electronics, optics, and mechanical components. Shock tests simulate these environments using electrodynamic shakers or actual pyrotechnic devices on test fixtures.}{catSS}

\dictentry{Sidereal Period}{The time required for a celestial body to complete one orbit relative to the distant fixed stars, as opposed to the synodic period measured relative to the Sun. Earth's sidereal period is approximately 365.256 days, while the Moon's sidereal period is about 27.3 days. The sidereal period is the true orbital period and is used as the baseline for calculating other orbital timing relationships.}{catSS}

\dictentry{Slingshot}{An informal term for a gravity assist maneuver in which a spacecraft uses the gravitational field and orbital motion of a planet to change its velocity relative to the Sun without expending propellant. The spacecraft follows a hyperbolic trajectory past the planet, gaining or losing heliocentric speed depending on the approach geometry. The Voyager missions famously used a series of slingshot maneuvers around the outer planets to visit multiple targets on a single trajectory.}{catSS}

\dictentry{SmallSat}{A broad category of small satellites encompassing miniaturized spacecraft with mass typically under 500 kilograms, including microsats, nanosats, and cubesats. SmallSats leverage miniaturized electronics, commercial off-the-shelf components, and reduced-size subsystems to dramatically lower development cost and time. The rapid growth of smallsat constellations for internet, Earth observation, and IoT applications has fundamentally changed the economics and accessibility of space.}{catSS}

\dictentry{Soft Capture}{The initial phase of a docking sequence in which the active spacecraft's capture mechanism gently engages with the passive spacecraft's docking interface with minimal relative velocity and force. Soft capture typically involves a probe, latches, or magnetic guides that absorb residual approach energy and establish a preliminary connection before the structural hard-mate sequence begins. This gentle contact is critical for preventing damage to delicate docking mechanisms and sensitive payloads.}{catSS}

\dictentry{Solar Array}{An assembly of interconnected solar panels and support structure on a spacecraft that converts sunlight into electrical power to operate onboard systems and charge batteries. Solar arrays may be body-mounted for simplicity or deployable wing-type structures for higher power generation, and must be sun-pointed for maximum efficiency. The total power output depends on cell efficiency, array area, solar distance, sun angle, temperature, and degradation from radiation exposure over the mission lifetime.}{catSS}

\dictentry{Solar Array Deployment}{The sequence of mechanically unfurling or extending a spacecraft's stowed solar panels to their full operational configuration after launch, typically one of the most critical and high-risk events in a mission. Solar arrays are folded to fit within the launch vehicle fairing and must deploy reliably using a combination of spring hinges, motorized drives, and hold-down release mechanisms. Deployment is usually commanded soon after spacecraft separation and verified through telemetry and visual confirmation.}{catSS}

\dictentry{Solar Flare}{A sudden, intense brightening on the Sun's surface caused by the release of magnetic energy stored in the solar atmosphere, emitting radiation across the entire electromagnetic spectrum. Solar flares accelerate particles to near-relativistic speeds and can last from minutes to hours, producing X-ray and ultraviolet radiation that reaches Earth in eight minutes and can disrupt radio communications. Major flares are often associated with coronal mass ejections that cause geomagnetic storms days later.}{catSS}

\dictentry{Solar Panel}{A flat assembly of photovoltaic cells that directly converts sunlight into electricity through the photovoltaic effect, forming the power source for most spacecraft in the inner solar system. Space-rated solar panels use multi-junction gallium arsenide cells achieving efficiencies above 30 percent, protected by cover glass against radiation damage and micrometeorite impacts. The panels must be oriented toward the Sun and their power output varies with the square of the distance from the Sun.}{catSS}

\dictentry{Solar Sail}{A spacecraft propulsion method that uses the pressure of sunlight on a large, thin, reflective membrane to generate continuous thrust without propellant. The sail acts as a mirror, reflecting photons and transferring their momentum to the spacecraft, producing a small but constant acceleration that can build to high velocities over time. Solar sail technology has been demonstrated by missions such as JAXA's IKAROS and NASA's NanoSail-D and holds promise for long-duration interstellar precursor missions.}{catSS}

\dictentry{Solar System}{The gravitationally bound system of the Sun and all objects that orbit it, including eight planets, their moons, dwarf planets, asteroids, comets, and interplanetary dust, spanning from the Sun to the Oort Cloud at the outermost boundary. The inner solar system contains the rocky terrestrial planets Mercury, Venus, Earth, and Mars, while the outer solar system hosts the gas and ice giants Jupiter, Saturn, Uranus, and Neptune. The solar system formed approximately 4.6 billion years ago from the gravitational collapse of a molecular cloud.}{catSS}

\dictentry{Solar Wind}{A continuous stream of charged particles, primarily protons and electrons, flowing outward from the Sun's corona at speeds of 300 to 800 kilometers per second. The solar wind fills the heliosphere, carries the Sun's magnetic field into interplanetary space, and interacts with planetary magnetospheres to create phenomena such as auroras and geomagnetic storms. Variations in solar wind density and speed with the solar cycle significantly affect space weather and satellite operations.}{catSS}

\dictentry{Southern Lights}{The Aurora Australis, the southern hemisphere counterpart of the Aurora Borealis, caused by charged particles from the solar wind being funneled along Earth's magnetic field lines into the southern polar atmosphere. The displays produce similar colors and patterns as the northern aurora, though they are less frequently observed due to the remoteness of Antarctic viewing locations. Simultaneous observations of both auroras provide valuable data about the global structure of Earth's magnetosphere.}{catSS}

\dictentry{Space Probe}{An unmanned spacecraft that is sent on a trajectory to explore celestial bodies beyond Earth orbit, typically performing flybys, orbital insertion, or atmospheric entry at one or more targets. Space probes carry scientific instruments and communication equipment to transmit findings back to Earth, often operating autonomously due to the long communication delays involved. Historic probes include Voyager 1 and 2, which have traveled beyond the heliopause into interstellar space.}{catSS}

\dictentry{Space Science}{The interdisciplinary scientific study of the space environment, including the physics of the Sun, heliosphere, magnetospheres, upper atmospheres, and the nature of celestial bodies throughout the universe. Space science encompasses solar physics, space plasma physics, space astronomy, and planetary science, advanced through both robotic missions and crewed exploration. Findings from space science have transformed our understanding of the cosmos and continue to address fundamental questions about the universe, from black holes to the origins of planetary systems.}{catSS}

\dictentry{Space Station}{A large, habitable artificial satellite in orbit designed for long-duration human presence and the conduct of scientific experiments in microgravity and the space environment. Space stations are assembled in orbit from multiple pressurized modules and are serviced by regular crew rotation and cargo resupply missions. The International Space Station, continuously crewed since 2000, represents the largest collaborative engineering project in human history.}{catSS}

\dictentry{Space Suit}{A full-body pressurized garment that provides life support, thermal control, micrometeorite protection, and mobility for astronauts performing extravehicular activities outside a spacecraft. The suit maintains an internal atmosphere, provides oxygen and removes carbon dioxide, manages thermal loads through a liquid cooling garment and sublimator, and shields against radiation and small particle impacts. Modern EVA suits weigh over 100 kilograms on Earth but allow astronauts to work productively for up to eight hours in the vacuum of space.}{catSS}

\dictentry{Space Tether}{A long, thin cable or ribbon extending between two or more spacecraft or from a spacecraft to a counterweight, used for a variety of purposes including artificial gravity generation, electrodynamic propulsion, and momentum exchange. Tethers can be tens to hundreds of kilometers long and must withstand the tension forces and space environment including atomic oxygen erosion and micrometeorite impacts. The concept was demonstrated by the TSS-1 and TSS-1R missions, though practical challenges remain.}{catSS}

\dictentry{Space Tourism}{The emerging industry of providing commercial spaceflight experiences to private individuals for recreational, leisure, or educational purposes, moving beyond government astronaut programs. Companies such as SpaceX, Blue Origin, and Virgin Galactic are developing vehicles capable of carrying paying passengers on suborbital and orbital flights. Space tourism is driving down launch costs and advancing vehicle reusability, potentially making space access available to a much broader segment of the population.}{catSS}

\dictentry{Spacecraft}{A vehicle or machine designed to travel and operate in the vacuum of space, consisting of integrated subsystems for power, propulsion, thermal control, communication, navigation, and payload. Spacecraft range from tiny cubesats weighing a few kilograms to massive interplanetary probes, and may be crewed or uncrewed. Each subsystem must function reliably in the harsh space environment of vacuum, radiation, extreme temperatures, and microgravity for the duration of the mission.}{catSS}

\dictentry{Spacecraft Bus}{The primary structural framework and core subsystems of a spacecraft that provide the necessary support functions for the payload or mission instruments, including power, thermal control, attitude control, propulsion, and communications. The bus is the platform upon which mission-specific payloads are mounted and integrated, and its design is often reused across multiple missions with different payloads to reduce development cost and risk.}{catSS}

\dictentry{Spaceport}{A specialized ground facility designed for launching, recovering, and servicing spacecraft and launch vehicles, analogous to an airport for aircraft. Spaceports include launch pads, processing facilities, mission control centers, and associated infrastructure such as propellant storage and tracking networks. Modern commercial spaceports are being developed worldwide to support the growing demand for launch services from both government and commercial operators.}{catSS}

\dictentry{Sphere of Influence}{The region around a celestial body within which its gravitational influence dominates over that of other bodies for the purposes of trajectory analysis. Within a body's sphere of influence, its gravity is treated as the primary force, and beyond it, the gravitational pull of the central body takes over. Earth's sphere of influence extends to about 925,000 kilometers, while Jupiter's extends over 50 million kilometers, reflecting its enormous mass.}{catSS}

\dictentry{Spin Balance}{The process of verifying and adjusting the mass distribution of a spinning spacecraft to ensure that its principal spin axis coincides with the geometric axis, eliminating wobble. An unbalanced spinning spacecraft will precess like a poorly thrown football, causing pointing errors and structural stresses. Spin balance is achieved through iterative adjustment of small balance masses until the measured wobble falls within specified tolerances, using precision spin tables and vibration analysis equipment.}{catSS}

\dictentry{Spin Stabilization}{A method of maintaining spacecraft attitude by spinning the vehicle about one axis, exploiting the gyroscopic rigidity that resists changes to the spin axis direction. Spin-stabilized spacecraft maintain a fixed orientation in inertial space without requiring active control systems or propellant for most attitude maintenance. This approach was used extensively on early communications satellites and remains useful for simple missions, though it limits the ability to point instruments independently.}{catSS}

\dictentry{Spin Table}{A precision rotating platform used on the ground to spin up spacecraft and their payloads to operational rotation rates for balance verification and deployment testing. The spin table allows engineers to measure mass balance, verify deployment mechanisms under centrifugal loading, and confirm that spinning appendages deploy correctly at the expected rates. It is an essential piece of ground support equipment for all spin-stabilized missions and also supports dynamic balancing procedures.}{catSS}

\dictentry{Star}{A luminous sphere of plasma held together by its own gravity, generating energy through nuclear fusion reactions that convert hydrogen into helium in its core. Stars range in mass from about 0.08 to over 100 times the Sun's mass, with their luminosity, temperature, color, and lifetime determined primarily by their initial mass. Stars are the fundamental building blocks of galaxies and the source of the heavy elements necessary for planet formation and life.}{catSS}

\dictentry{Star Tracker}{A highly sensitive optical attitude determination sensor that photographs star fields and compares observed star patterns against an onboard star catalog to determine the spacecraft's absolute orientation. Star trackers achieve accuracy of arcseconds or better and are the most precise attitude sensors commonly used on modern spacecraft. They operate autonomously, typically identifying several dozen stars in each image and computing a three-axis attitude solution within seconds.}{catSS}

\dictentry{Stratosphere}{The layer of Earth's atmosphere extending from about 12 to 50 kilometers altitude, above the troposphere and below the mesosphere. The stratosphere contains the ozone layer, which absorbs ultraviolet radiation from the Sun and causes temperature to increase with altitude. Commercial aircraft occasionally fly in the lower stratosphere to avoid weather, and weather balloons and some research aircraft reach the upper stratosphere.}{catSS}

\dictentry{Structural Model}{A full-scale spacecraft structural representation built for static and dynamic load testing to verify that the structure can withstand the forces and vibrations experienced during launch, ascent, and orbital maneuvers. The structural model may be built with flight-representative materials and joint details but does not include flight electronics or mechanisms. Test results validate analytical structural models and confirm that design margins of safety are adequate.}{catSS}

\dictentry{Structure}{The mechanical framework of a spacecraft that provides the primary load paths for launch and operational forces, supports all subsystems and payloads, and maintains the geometric relationships between components. Spacecraft structures are typically made from aluminum alloys, titanium, honeycomb sandwich panels, or composite materials chosen for their stiffness-to-weight ratio, thermal stability, and radiation tolerance. Structural design must minimize mass while meeting strength, stiffness, and dynamic requirements.}{catSS}

\dictentry{Sun}{The star at the center of our solar system, a G-type main-sequence star containing 99.86 percent of the total solar system mass, with a diameter of about 1.39 million kilometers and a surface temperature of approximately 5,778 Kelvin. The Sun generates energy through hydrogen fusion in its core at a rate of about 3.8 times 10 to the 26 watts, producing the light and heat that sustain life on Earth. Its activity cycles through an approximately 11-year sunspot cycle that affects space weather throughout the solar system.}{catSS}

\dictentry{Sun Sensor}{A simple attitude determination device that detects the direction to the Sun using photodetectors or coded apertures, providing one-axis or two-axis orientation information relative to the Sun vector. Sun sensors are among the most common and reliable attitude sensors, used for initial attitude acquisition, safe-mode pointing, and solar array orientation. They range from simple analog devices measuring one axis to digital sensors with multiple pixels for coarse two-axis determination.}{catSS}

\dictentry{Sun-Synchronous Orbit}{A near-polar orbit in which the orbital plane precesses at the same rate as the Earth orbits the Sun, maintaining a constant angle between the orbit plane and the Sun direction. This geometry ensures that the satellite passes over any given latitude at the same local solar time on each orbit, providing consistent lighting conditions for Earth observation imagery. Sun-synchronous orbits are achieved at altitudes near 600 to 800 kilometers with inclinations of about 98 degrees.}{catSS}

\dictentry{Supernova}{A catastrophic stellar explosion that occurs at the end of a massive star's life or in certain binary star systems, releasing enormous amounts of energy that can briefly outshine an entire galaxy. Supernovae are classified as Type Ia, involving the thermonuclear detonation of a white dwarf, or Type II, involving the core collapse of a massive star. The explosion disperses heavy elements into the interstellar medium, enriching the material from which new stars and planets form.}{catSS}

\dictentry{Swing-By}{Another term for a gravity assist or gravitational slingshot maneuver, where a spacecraft uses the gravity and orbital velocity of a celestial body to change its trajectory and speed without using propellant. During a swing-by the spacecraft approaches the body, is deflected by its gravity, and departs with a changed velocity vector relative to the Sun. The Voyager mission's Grand Tour of the outer planets relied on a carefully choreographed sequence of planetary swing-bys to visit Jupiter, Saturn, Uranus, and Neptune.}{catSS}

\dictentry{Synchronous Rotation}{A condition in which a celestial body rotates on its axis in exactly the same time it takes to complete one orbit around its primary, causing the same face to always point toward the primary. Earth's Moon is tidally locked in synchronous rotation, which is why only the near side is visible from Earth without a spacecraft. This phenomenon results from tidal forces that gradually slow the body's rotation until it matches its orbital period.}{catSS}

\dictentry{Synodic Period}{The time interval between two successive identical configurations of a celestial body as seen from the Sun, such as two successive oppositions or two successive conjunctions. The synodic period differs from the sidereal period because the observer's reference point is also moving. For Mars, the synodic period is about 780 days, defining the time between successive Mars oppositions when the planet is closest to Earth and best positioned for observation.}{catSS}

\dictentry{Telecommand}{Encoded commands transmitted from a ground station to a spacecraft to control its functions, configure subsystems, execute maneuvers, and respond to anomalies. Telecommands must be received, verified for integrity, and executed by the spacecraft's onboard computer, often with redundancy checks to prevent erroneous operations. Due to light-speed delays, deep-space telecommands must be sent hours before the desired execution time.}{catSS}

\dictentry{Telemetry}{Data collected by onboard sensors and instruments on a spacecraft and transmitted to ground stations for monitoring, analysis, and record-keeping. Telemetry includes measurements of temperatures, voltages, pressures, attitude, orbital parameters, and payload data, forming the primary means by which ground controllers assess spacecraft health and performance. Telemetry rates vary from a few bits per second for simple spacecraft to megabits per second for data-intensive science missions.}{catSS}

\dictentry{Terrestrial Planet}{A rocky planet composed primarily of silicate rocks and metals, with a solid surface and relatively high density, exemplified by Mercury, Venus, Earth, and Mars in our solar system. Terrestrial planets formed in the inner solar system where temperatures were too high for volatile ices to condense, resulting in compositions dominated by refractory materials. Their relatively small sizes and masses compared to gas giants are a consequence of the limited material available in the inner protoplanetary disk.}{catSS}

\dictentry{Test}{The systematic process of verifying that spacecraft hardware, software, and systems meet their design requirements and can survive and operate in the intended space environment. Testing progresses from component-level screening through subsystem verification to full spacecraft environmental testing, including vibration, acoustic, thermal vacuum, and electromagnetic compatibility tests. A rigorous test campaign is essential for identifying defects and validating the flight design before committing to launch.}{catSS}

\dictentry{Thermal Blanket}{An insulating cover made from multi-layer insulation or other thermal control materials that is wrapped around spacecraft components to control heat transfer between the component and its environment. Thermal blankets reduce radiative heat loss in cold conditions or prevent excessive solar heating, maintaining components within their operational temperature range. They are among the most commonly used passive thermal control elements on all types of spacecraft and are often custom-fabricated to fit complex geometries.}{catSS}

\dictentry{Thermal Control System}{An integrated set of hardware that maintains all spacecraft components within their required temperature ranges by managing heat generation, transfer, and rejection. The system includes passive elements such as insulation and coatings, active elements such as heaters and cold plates, and heat transport mechanisms including heat pipes and fluid loops. Effective thermal control is critical because electronic, mechanical, and optical components each have specific temperature requirements for proper operation.}{catSS}

\dictentry{Thermal Model}{A mathematical representation of a spacecraft's thermal behavior that predicts temperature distributions under various operating conditions and environmental scenarios. Thermal models use node-based network analysis to simulate heat transfer through conduction, radiation, and convection, and are validated against thermal vacuum test data. They are essential tools for designing the thermal control system and predicting temperatures throughout the mission lifetime.}{catSS}

\dictentry{Thermal Vacuum Test}{An environmental test in which a spacecraft is placed in a vacuum chamber and subjected to the extreme temperature cycles of the space environment, with hot and cold walls simulating solar heating and deep-space cold. The test verifies that all components operate correctly across their full temperature range and that thermal control hardware functions as designed. Thermal vacuum testing is typically the longest and most resource-intensive test in the spacecraft qualification campaign.}{catSS}

\dictentry{Thermosphere}{The layer of Earth's atmosphere extending from about 85 to 600 kilometers altitude, above the mesosphere, where temperatures rise dramatically due to absorption of extreme ultraviolet radiation from the Sun. Despite the high kinetic temperatures, the thermosphere is extremely thin and would feel very cold to a human because of the low particle density. The International Space Station orbits within the lower thermosphere, where residual atmospheric drag gradually decays its orbit.}{catSS}

\dictentry{Third Body Perturbation}{An orbital disturbance caused by the gravitational influence of a celestial body other than the primary body around which the satellite orbits, such as the Sun's effect on an Earth satellite or Jupiter's effect on an asteroid. Third-body perturbations can cause long-term secular changes in orbital elements as well as periodic oscillations, and must be modeled for accurate long-term orbit prediction. For geostationary satellites, solar and lunar third-body effects are among the dominant perturbation sources.}{catSS}

\dictentry{Three-Axis Stabilization}{An active attitude control method that maintains a spacecraft's orientation in all three rotational axes using sensors and actuators such as reaction wheels, control moment gyros, and thrusters. Unlike spin stabilization, three-axis control allows the spacecraft body to remain non-spinning, enabling precise pointing of instruments and antennas in any direction. This method is used on most modern spacecraft because it provides the flexibility needed for complex mission operations.}{catSS}

\dictentry{Thruster}{A small rocket engine used for spacecraft attitude control, translation, orbit maintenance, and momentum desaturation, typically producing thrust levels from millinewtons to hundreds of newtons. Thrusters use various propellant types including cold gas, monopropellant hydrazine, or bipropellant combinations depending on performance, lifetime, and cost requirements. Electric propulsion thrusters such as ion engines and Hall-effect thrusters provide much higher specific impulse for fuel-efficient long-duration missions.}{catSS}

\dictentry{Tidal Force}{A differential gravitational force that stretches an object along the direction toward the source of gravity and compresses it perpendicular to that direction, caused by the variation in gravitational pull across the object's extent. Tidal forces are responsible for the ocean tides on Earth, the volcanic activity on Jupiter's moon Io, and the tidal heating of Europa's subsurface ocean. In extreme cases near massive objects, tidal forces can tear apart bodies that cross the Roche limit.}{catSS}

\dictentry{Tidal Locking}{The process by which tidal forces gradually synchronize a satellite's rotation period with its orbital period, causing the same hemisphere to permanently face the primary body. Most large moons in the solar system are tidally locked, including Earth's Moon, which is why we only see its near side from Earth. Tidal locking transfers angular momentum between the orbit and the rotation over millions to billions of years until equilibrium is reached.}{catSS}

\dictentry{Transponder}{A combined transmitter and receiver unit on a communications satellite that receives uplinked signals from a ground station, amplifies them, shifts their frequency, and retransmits them back to Earth on a different frequency. Transponders are the fundamental channel units of a communications satellite, with each transponder handling a specific bandwidth of traffic. A typical communications satellite carries dozens of transponders, enabling thousands of simultaneous voice, data, and video channels.}{catSS}

\dictentry{Troposphere}{The lowest layer of Earth's atmosphere, extending from the surface to about 12 kilometers altitude at mid-latitudes, where nearly all weather phenomena occur. Temperature decreases with altitude in the troposphere due to heating of the surface by solar radiation, creating the convection currents that drive cloud formation, precipitation, and storm systems. The troposphere contains approximately 75 percent of the total atmospheric mass and nearly all of its water vapor.}{catSS}

\dictentry{True Anomaly}{The angular position of a body in its orbit measured from the periapsis point in the direction of motion, representing the actual position of the satellite along its elliptical path. Unlike mean anomaly, true anomaly accounts for the non-uniform speed of a body in an elliptical orbit, increasing rapidly near periapsis and slowly near apoapsis. Converting between true anomaly and other anomaly types is a fundamental calculation in orbital mechanics.}{catSS}

\dictentry{Tundra Orbit}{A highly elliptical geosynchronous orbit with an inclination of about 63.4 degrees that provides long dwell times over high-latitude regions, similar to the Molniya orbit but with a 24-hour period matching Earth's rotation. The tundra orbit's frozen argument of perigee keeps the apogee fixed over the northern hemisphere, allowing the satellite to spend the majority of each orbit visible from high-latitude ground stations. It is used for communications and broadcasting to polar regions where GEO satellites have poor visibility.}{catSS}

\dictentry{Turbo Code}{A class of high-performance error-correcting codes that achieve performance very close to the theoretical Shannon limit through the parallel concatenation of two recursive systematic convolutional codes with an interleaver. Turbo codes are decoded using iterative soft-input soft-output algorithms that exchange information between the two component decoders. They are used in deep-space communication standards and in cellular wireless standards such as LTE for their near-capacity performance.}{catSS}

\dictentry{Vacuum}{A space environment characterized by the near-complete absence of matter, with pressure far below atmospheric pressure at Earth's surface. The vacuum of space presents unique challenges for spacecraft including outgassing of materials, electrical charging, thermal management through radiation only, and the inability to use convective cooling. Even the relatively low orbit of the ISS maintains an extremely thin atmosphere that causes measurable drag on the station over time.}{catSS}

\dictentry{Van Allen Belt}{One of two or more torus-shaped zones of energetic charged particles trapped by Earth's magnetic field, consisting of an inner belt of high-energy protons and electrons at about 1,000 to 6,000 kilometers altitude and an outer belt of electrons extending from about 13,000 to 60,000 kilometers. The belts were discovered by James Van Allen in 1958 and pose radiation hazards to spacecraft electronics and astronauts. Spacecraft must be radiation-hardened to survive passage through and operation within the belts.}{catSS}

\dictentry{Vandenberg SFB}{Vandenberg Space Force Base, a United States Space Force launch facility on the central coast of California used for polar and retrograde orbit launches where the trajectory carries the vehicle southward over the Pacific Ocean. Its location at approximately 34.7 degrees north latitude makes it unsuitable for equatorial orbits but ideal for Sun-synchronous and reconnaissance satellite missions. Vandenberg has been a primary West Coast launch site since the early days of the space age.}{catSS}

\dictentry{Vibration Test}{An environmental test that subjects spacecraft hardware to the random and sinusoidal vibration environments expected during rocket launch, typically using electrodynamic shakers to reproduce the vibration spectra measured on previous flights. Vibration testing verifies that all components, structures, and mechanisms can withstand the intense dynamic loads without structural failure or functional degradation. The test levels are derived from launch vehicle acoustic and structural analyses and are applied in three orthogonal axes.}{catSS}

\dictentry{Wallops Flight Facility}{A NASA rocket launch and research facility located on Wallops Island on the Eastern Shore of Virginia, used for suborbital rocket launches, small orbital missions, and aeronautical research. WFF is one of the oldest launch sites in the United States, operational since 1945, and supports scientific sounding rockets, balloon launches, and commercial cargo resupply missions to the ISS. Its location provides launch trajectories over the Atlantic Ocean for safety.}{catSS}

\dictentry{Weightlessness}{The sensation or condition of apparent zero gravity experienced by objects in free fall, caused by the fact that everything in the local frame falls at the same rate regardless of mass. Astronauts on the ISS experience continuous weightlessness because the station and everything inside it are in constant free fall around Earth. True weightlessness does not exist near massive bodies, as residual accelerations from atmospheric drag, tidal effects, and crew activity create a microgravity environment rather than perfect zero-g.}{catSS}

\dictentry{White Dwarf}{The dense stellar remnant left behind when a low- to medium-mass star exhausts its nuclear fuel and sheds its outer layers as a planetary nebula, leaving behind a hot, compact core composed primarily of carbon and oxygen. White dwarfs are roughly Earth-sized but contain about half the Sun's mass, resulting in extraordinary densities where a teaspoon of material weighs several tons. They slowly cool and fade over billions of years, supported against gravitational collapse by electron degeneracy pressure.}{catSS}

\dictentry{X Band}{A portion of the electromagnetic spectrum in the frequency range from approximately 8 to 12 gigahertz, used for high-data-rate spacecraft communication and radar applications. X-band offers a good compromise between the high data rates of Ka-band and the weather resilience of lower-frequency bands, making it widely used for deep-space links and military radar. X-band frequencies are also used for planetary radar mapping and spacecraft tracking.}{catSS}

\dictentry{Yarkovsky Effect}{A subtle force on small rotating asteroids caused by the anisotropic emission of thermal radiation, where the afternoon side of the asteroid is warmer than the morning side and radiates more momentum. Though tiny, this effect accumulates over millions of years to significantly alter the orbits of small asteroids, playing an important role in delivering near-Earth asteroids from the main belt. It was first confirmed through precise radar tracking of asteroid 6489 Golevka.}{catSS}

\dictentry{Yo-Yo Deploy}{A spin-despin technique in which masses on cables are released from a spinning spacecraft and move outward due to centrifugal force, extracting angular momentum and reducing the spacecraft's spin rate before the cables and masses are jettisoned. The technique is simple, reliable, and requires no power or pyrotechnics, making it ideal for separating spin-stabilized upper stages from their payloads. The final spin rate depends on the cable length, mass, and the number of wire pairs used.}{catSS}

